\documentclass[12pt]{article}

\usepackage{amsmath}
\usepackage{lipsum}
\usepackage{tikz-cd} 


\usepackage{amsthm}
\usepackage{mathrsfs}	%for \mathscr{} script
\usepackage{amssymb}
\usepackage{amscd}
\usepackage{setspace}	%Pro­vides sup­port for set­ting the spac­ing be­tween lines in a doc­u­ment with the \sin­glespac­ing, \one­half­s­pac­ing, and \dou­blespac­ing com­mands. Other size spac­ings also avail­able.
\usepackage{tikz}		%graph drawing
\usepackage{tikz-cd}
\usetikzlibrary{matrix}	%diagram drawing using tikz
\usepackage{graphicx}	%insert images
\graphicspath{ {pics/} }	%local image folder
\usepackage{wrapfig}	%figure wrapping for text around pictures
\usepackage{enumerate}	%to get enumerate to use roman numerals [i], [I], [i.], [i)], [(i)]
\usepackage{bbm}		%to get blackboard 1 \mathbbm{1}
%\usepackage{hyperref}

\usepackage[utf8]{inputenc}	% for chapters, theorems and proofs
\usepackage[english]{babel}	% for theorems
\usepackage{blindtext}		% for chapters
\usepackage[affil-it]{authblk}	%for affiliation

\usepackage{faktor}			%for modding out stuff command \faktor{n}{d} left
\usepackage{xparse}			%for the macro \rfactor below right

\usepackage{hyperref}		%for adding hyperlinks
%\usepackage{cite}			%for bibliography
\usepackage[lite]{amsrefs}


\newtheorem{theorem}{Theorem}[section] %defines new theorem environment called with \theorem{} and displays Theorem

\newtheorem{definition}[theorem]{Definition}
%\numberwithin{definition}{subsection}

\newtheorem{lemma}[theorem]{Lemma}
%\numberwithin{lemma}{subsection}

\newtheorem{remark}[theorem]{Remark}
%\numberwithin{remark}{subsection}

\newtheorem{proposition}[theorem]{Proposition}
%\numberwithin{proposition}{subsection}

\newtheorem{corollary}[theorem]{Corollary}
%\numberwithin{corollary}{subsection}

\newtheorem{example}[theorem]{Example}

\newtheorem{examples}[theorem]{Examples}

\newtheorem{notation}[theorem]{Notation}

\newcommand{\acts}{\curvearrowright}
\newcommand{\asdim}{\text{asdim}}
\newcommand{\diam}{\text{diam}}
\newcommand{\R}{\mathbb{R}}
\newcommand{\Z}{\mathbb{Z}}
\newcommand{\C}{\mathbb{C}}
\newcommand{\N}{\mathbb{N}}
\newcommand{\dad}{\text{DAD}}

\newtheorem{thmx}{Theorem} %lettered theorems
\renewcommand{\thethmx}{\Alph{thmx}}

\newtheorem*{theorem*}{Example}

%\newtheorem{examplex}{Example} %lettered theorems
%\renewcommand{\thethmx}{\Alph{thmx}}

\DeclareFontFamily{U}{mathb}{\hyphenchar\font45}
\DeclareFontShape{U}{mathb}{m}{n}{ <-6> matha5 <6-7> matha6 <7-8>
mathb7 <8-9> mathb8 <9-10> mathb9 <10-12> mathb10 <12-> mathb12 }{}
\DeclareSymbolFont{mathb}{U}{mathb}{m}{n}

\DeclareMathAccent{\abxring}{0}{mathb}{"38}

\DeclareFontFamily{U}{mathb}{\hyphenchar\font45}
\DeclareFontShape{U}{mathb}{m}{n}{ <-6> matha5 <6-7> matha6 <7-8>
mathb7 <8-9> mathb8 <9-10> mathb9 <10-12> mathb10 <12-> mathb12 }{}
\DeclareSymbolFont{mathb}{U}{mathb}{m}{n}

\renewcommand{\thethmx}{\Alph{thmx}}

\begin{document}
\title{A Hurewicz-type Theorem for the Dynamic Asymptotic Dimension with Applications to Coarse Geometry and Dynamics}
\author{Samantha Pilgrim}
\maketitle

\paragraph{Abstract:} We prove a Hurewicz-type theorem for the dynamic asymptotic dimension originally introduced by Guentner, Willett, and Yu.  Calculations of (or simply upper bounds on) this dimension are known to have implications related to cohomology of group actions and the $K$-theory of their transformation group $C^*$-algebras.  Moreover, these implications are relevant to the current classification program for $C^*$-algebras.  As a corollary of our main theorem, we show the dynamic asymptotic dimension of actions by groups on profinite completions along sequential filtrations by normal subgroups is subadditive over extensions of groups, which shows that many such actions by elementary amenable groups are finite dimensional.  We combine this with other novel results relating the dynamic asymptotic dimension of such an action to the asymptotic dimension of a corresponding box space.  This allows us to give upper bounds on the asymptotic dimension of many box spaces (including examples from infinitely-many groups with exponential growth) using a generalization of the Hirsch length for elementary amenable groups.  For some of these examples, we can also find lower bounds by utilizing the theory of ends of groups.  

\section{Introduction}
When considering any invariant thought to behave like dimension, a natural question is whether it is subadditive over morphisms.  That is, given a morphism $f: X\to Y$ of two objects, is the dimension of $X$ at most the dimension of $Y$ plus the largest dimension of a `fibre' of $f$  (speaking loosely)?  An affirmative answer to such a question is sometimes called a Hurewicz-type theorem, and speaks to the naturality and dimension-theoretic nature of such an invariant.  Classical results of this kind include the rank-nullity theorem (for the category of vector spaces and linear maps) and the Hurewicz mapping theorem \cite[Theorem 1.12.4, page 136]{engelking} (for the category of topological spaces and continuous maps).  Special cases include whether a notion of dimension is subadditive over products or extensions of objects.  

More recently, Hurewicz-type theorems have been developed for Gromov's asymptotic dimension in \cite{asymptotichurewicztheorem} and later in \cite{Brodskiy2006AHT}.  This paper uses similar methods as in \cite{Brodskiy2006AHT}, together with some original ideas, to prove a Hurewicz-type theorem for the dynamic asymptotic dimension ($\dad$) of Guentner, Willett, and Yu.  This dimension theory was introduced in \cite{dasdimGWY} as a dynamical analogue of the asymptotic dimension.  It is related to conditions used in work on the Farrell-Jones conjecture on manifold topology \cite{bartels} \cite{bartels_et_al}, as well as to the nuclear dimension of transformation group $C^*$-algebras \cite[Theorem 8.6]{dasdimGWY}.  

We state the appropriate Hurewicz-type theorem for the dynamic asymptotic dimension below.  The freeness assumption is negligible, as we explain in \ref{freeness remark}.  

\begin{thmx}
Suppose $\Gamma\acts X$ and $\Lambda\acts Y$ are actions on compact metric spaces by homeomorphisms which are each either free or isometric, and that $f: X\to Y$ and $\alpha: \Gamma\times X\to \Lambda$ are continuous functions such that $f(\gamma\cdot x) = \alpha(\gamma, x)\cdot f(x)$ for all $x\in X$ and $\gamma\in \Gamma$.  If $\dad(f):=\sup_{A\subset Y : \dad(A) = 0} \dad_{free}(f^{-1}(A))$, it holds that 

$$\dad_{free}(\Gamma\acts X)\leq \dad(f) + \dad_{free}(\Lambda\acts Y).\footnote{The notation $\dad_{free}$ refers to the definition of $\dad$ for free actions, see \ref{dad definition}.  }$$  
\end{thmx}\label{Theorem A}

Apart from being a natural permanence property for the dynamic asymptotic dimension, this theorem is interesting because it provides a new tool for showing certain actions have finite $\dad$.  While it is thought that an action by an amenable group with finite asymptotic dimension should have finite $\dad$, relatively little progress has been made aside from the case where the acting group is virtually nilpotent \cite[Theorem 7.4]{szabo2017rokhlin}.  Finiteness of the dynamic asymptotic dimension is important to the current classification program for $C^*$-algebras because of how $\dad$ relates to nuclear dimension and computations of $K$-theory via the Baum-Connes and HK conjectures \cite{Guentner2016DynamicAD} \cite{bonicke2021dynamic}.  New examples of finite dimensional actions add to the class of algebras which are both classifiable by their $K$-theoretic and tracial data and which admit extensive technology for computing $K$-theory.  $\dad$ also plays a special role for those approaching classification from a dynamical perspective, as it applies to actions with finite stabilizers which are not essentially free, and therefore not covered by Kerr's notion of almost finitness (see for instance \cite{MR4066584} and \cite{https://doi.org/10.48550/arxiv.2107.05273}).  The introduction of \cite{https://doi.org/10.48550/arxiv.2201.03409} gives a precise statement of the classification theorem and additional references.  %HK conjecture

We also give two applications of this new Hurewicz-type theorem.  First, we can show that the $\dad$ is preserved when one restricts to a finite-index subgroup of the acting group, similar to how the asymptotic dimension of a finite index subgroup is the same as that of the larger group.  This further demonstrates that the $\dad$ is a coarse dimension theory for dynamics.  

\begin{thmx} Suppose $\Gamma\acts X$ is a free action on a compact metric space by homeomorphisms and $\Lambda<\Gamma$ is a finite index subgroup.  Then $\dad(\Gamma\acts X)= \dad(\Lambda\acts X)$.  
\end{thmx}\label{Theorem B}

\noindent We can also show that the $\dad$ of groups acting on their profinite completions (sometimes called odometer actions) behaves subadditively when one has an extension of groups:

\begin{thmx}
Suppose $1\to \Delta\to \Gamma\to \Lambda\to 1$ is an exact sequence of countable groups and that $\Gamma$ is residually finite.  Suppose $(N_i)$ is a collection of finite index normal subgroups of $\Gamma$ directed by inclusion which is also a filtration of $\Gamma$.  Then $(N_i)$ induces collections of finite index subgroups and hence odometers of $\Delta$ and $\Lambda$.   If $\text{DAD}_{free}(\Delta\acts \widehat{\Delta}_{(N_i\cap \Delta)}) = m$ and $\text{DAD}_{free}(\Lambda\acts \widehat{\Lambda}_{(N_i/\Delta\cap N_i)}) = n$, $\text{DAD}_{free}(\Gamma\acts \widehat{\Gamma}_{(N_i)}) \leq m + n$.  
\end{thmx}\label{Theorem C}

This extension theorem can further be applied to coarse geometry, specifically the study of box spaces of groups.  These are spaces constructed out of finite quotients of a group.  Just as the asymptotic dimension of a group itself can have interesting group-theoretic applications (see for instance \cite{MR2431020}), there are likewise many correspondences between geometric properties of box spaces and group-theoretic properties.  For example, coarse equivalence of box spaces encodes virtual isomorphism of groups \cite{MR4091056}, and amenability of groups is equivalent to property A of their box spaces \cite[Propsotion 11.39]{Roe2003LecturesOC}.  Similar results also exist for a-T-menablity \cite{https://doi.org/10.1112/blms/bdt045}.  

Box spaces provided some of the first examples of spaces which do not coarsely embed in hilbert space \cite[Remark 11.27]{Roe2003LecturesOC} and it has been shown that a group with a finite dimensional box space must be amenable \cite[Remark 11.38 and Proposition 11.39]{Roe2003LecturesOC}.  Similar to how actions by finite dimensional amenable groups are thought to be finite dimensional, it is thought that a partial converse to this latter fact should also hold; that is, an amenable group with finite asymptotic dimension ought to have finite dimensional box spaces.  Notwithstanding, there has been similarly little progress on this aside from the case of virtually nilpotent groups \cite{boxspacesDT}.  

As it turns out, questions about the asymptotic dimension of box spaces are essentially special cases of questions about the dynamic asymptotic dimension, echoing a relationship between DAD and the asymptotic dimension of warped cones utilized in \cite[Section 8]{warpedcones}.  More precisely, we prove the following.  

\begin{thmx} \label{Theorem D}
Let $\Gamma$ be a countable group and $(N_i)$ a sequence of finite index subgroups which is directed by inclusion.  Suppose $\Gamma\acts \widehat{\Gamma}_{(N_i)}$ by left multiplication.  Then $\dad_{free}(\Gamma\acts \widehat{\Gamma}_{(N_i)}) = \asdim\Box_{(N_i)}\Gamma$.  
\end{thmx}

This relationship allows us to prove an extension theorem for the asymptotic dimension of box spaces, at least for box spaces along sequential filtrations by normal subgroups.  In fact, such an extension theorem was our original motivation for investigating a Hurewicz-type theorem for the DAD.  A proof of a similar result was attempted in \cite{finnsell2015asymptotic}, but the proof therein contains a gap which we discuss later in this section.  

\begin{thmx}
Suppose $1\to \Delta\to \Gamma\xrightarrow{\pi} \Lambda\to 1$ is an exact sequence of countable groups and $(N_i)$ is a sequence of finite index normal subgroups of $\Gamma$ directed by inclusion which is also a filtration of $\Gamma$.  This sequence induces sequences and hence box spaces of $\Delta$ and $\Lambda$.   Then if $\asdim\Box_{(N_i\cap \Delta)}\Delta = m$ and $\asdim\Box_{(\pi(N_i))}\Lambda = n$, $\asdim\Box_{(N_i)}\Gamma \leq m + n$.  
\end{thmx} 

We have therefore established a special case of the (unproved) statement \cite[Theorem A]{finnsell2015asymptotic}.  Our extra assumptions are that $\Gamma$ is countable and that $(N_i)$ is a sequence and a filtration of $\Gamma$ by normal subgroups.  In \cite[Theorem A]{finnsell2015asymptotic}, the family of subgroups need not be a sequence and are assumed only to be semiconjugacy separating.  This latter assumption is weaker than assuming $(N_i)$ to be a filtration by normal subgroups.  

Something should be said about how our approach differs from that of \cite{finnsell2015asymptotic}.  We employ a similar general strategy in attempting to apply a Hurewicz-type theorem, and as our notation will indicate, we are using similar arguments to handle DAD as are used in \cite{Brodskiy2006AHT} to handle asymptotic dimension.  Notwithstanding, our use of a dynamic Hurewicz theorem rather than a geometric one allows us to avoid the issue with \cite[Lemma 4.1(5)]{finnsell2015asymptotic}.  An example first provided by Jiawen Zhang and restated below shows the aforementioned statement is false.  We are very grateful to Jianchao Wu for helping to communicate this example and explain the problem that it highlights.  

\begin{theorem*}Let $6\Z$ and $9k\Z$ be subgroups of $\Z$ for $k$ odd.  Let $6\Z/(6\Z\cap 9k\Z)\to \Z/9k\Z$ be the natural inclusion (see the diagram in \cite[Remark 4.2]{finnsell2015asymptotic}).  When considered with the quotient metrics induced from a metric on $\Z$, this inclusion distorts the length of elements by an arbitrarily large amount for $k$ sufficiently large.  \end{theorem*}

\normalfont
The problem is essentially that, although an inclusion of a subgroup preserves asymptotic dimension, it does not do so quantitatively, that is, an inclusion of a subgroup is not in general an isometry.  Since the inclusions of a sequence of quotients need not preserve asymptotic dimension uniformly, any approach that tries to apply a Hurewicz-type theorem to a map of box spaces by analyzing the maps of the different components separately appears unworkable.  Indeed, one can attempt such an argument even using $\dad$ (as the author did previously) and will still encounter an analogous problem.  Our main theorem, however, can be applied more directly and so does not depend on uniform control over the maps of the individual components.  It appears, therefore, that there is some geometric obstruction which our dynamical approach circumvents.  

The extension theorem for the asymptotic dimension of box spaces provides new machinery for bounding the dimension of box spaces which applies to non virtually nilpotent groups.  Inspired by the techniques in \cite[Section 4]{finnsell2015asymptotic}, we use this machinery to derive several consequences for the asymptotic dimension of box spaces of elementary amenable groups.  We obtain a special case of \cite[Theorems B]{finnsell2015asymptotic} (a consequence of that paper's unproved main theorem) relating asymptotic dimension and Hirsch length:  

\begin{thmx} Let $\Gamma$ be a countable, residually finite, elementary amenable group.  Then for any sequence  $(N_i)$ of finite index normal subgroups which is also a filtration, $\asdim \Gamma\leq \asdim\Box_{(N_i)} \Gamma\leq h(\Gamma)$.  
\end{thmx}

The reverse inequality also holds in a more restrictive setting (a special case of \cite[Theorem C]{finnsell2015asymptotic}):

\begin{thmx}
If $\Gamma$ is countable and residually finite and $\Gamma = \Lambda\rtimes \Delta$ where $\Lambda$ is virtually polycyclic and $\Delta$ is locally finite, then for any sequence $(N_i)$ of finite index normal subgroups which is also a filtration, $\asdim\Gamma = \asdim\Box_{(N_i)}\Gamma = h(\Gamma)$.  
\end{thmx}

These last two results show that the asymptotic dimension and the Hirsch length, both natural ways of numerically capturing the rank or complexity of a group, are in fact closely related, building on previous work such as \cite[Theorem 3.5]{MR2213160}.  

We conclude by using the extension theorem in several example calculations.  We show that sequential box spaces of $BS(1, n)$ along filtrations have asymptotic dimension $2$ for all $n$, which in particular exhibits an infinite class of groups with exponential growth and finite-dimensional box spaces.  In obtaining the lower bound for this result, we compute the asymptotic dimension of $BS(1, n)$, which by itself involves some interesting methods in geometric group theory, specifically the notion of ends of groups.  We also deal with some examples of more complicated groups which have similar presentations or are in some sense `generalized Baumslag-Solitar groups'.  

\section{(Dynamic) Asymptotic Dimension}
The following concept is originally due to Gromov, specifically \cite[Section 1.E]{Gromov1991AsymptoticIO}.  There are many definitions, but we will use one throughout which makes the analogy with dynamic asymptotic dimension obvious:

\definition{Let $(X, d)$ be a metric space.  If $U\subset X$ is a subset, an $r$-chain in $U$ is a sequence $x_0, \ldots, x_n$ of points in $U$ such that $d(x_i, x_{i+1})\leq r$ for all $i$; and two points of $U$ are in the same $r$-component if they are connected by an $r$-chain in $U$.  The subset $U$ is said to be $r$-connected if any two points in $U$ are connected by an $r$-chain in $U$.  Finally, $\asdim X\leq n$ means for all $r>0$ there is $M>0$ and a cover $\mathcal{U} = \{U_0, \ldots, U_n\}$ of $X$ such that each $r$-component of each $U_i$ has diameter at most $M$.  Such a cover is called a $(d, R, M)$-cover of $X$.  The asymptotic dimension of $X$ is then the smallest integer $d$ such that $\asdim X\leq d$, and is infinite if there is no such $d$.  }

\notation{Throughout, $\Gamma\acts X$ and $\Lambda\acts Y$ will be actions on compact metric spaces by homeomorphisms.  }

\normalfont
We now define the dynamic asymptotic dimension of a group action, originally introduced by Guentner, Willett, and Yu in \cite{dasdimGWY}.  

\definition{Let $\Gamma\acts X$ be an action on a compact Hausdorff space by homeomorphisms.  The \textit{dynamic asymptotic dimension} $\Gamma\acts X$, denoted $\dad(\Gamma\acts X)$, is the smallest integer $d$ such that for all finite $F\subset \Gamma$ there exists an open cover $\mathcal{U} = \{U_0, \ldots, U_d\}$ of $X$ such that the set

$$
\Bigg\{\begin{array}{l|l} g \in \Gamma & g = f_k\cdots f_1 \text{ and } \exists x\in U_j\text{, } 0\leq j\leq d \\ & \text{ such that } f_{k_0}\cdots f_1\cdot x\in U_j \text{ }\forall \text{ } 1\leq k_0\leq k \end{array}\Bigg\}
$$

is finite.}

\normalfont
If $\Gamma\acts X$ is free, this definition is equivalent to the following: %Prove this!

\definition{If $\Gamma\acts X$ is an action on a compact space by homeomorphisms, $\dad_{free}(\Gamma\acts X)$ is the smallest integer $d$ such that for all finite $F\subset \Gamma$ there is an open cover $\mathcal{U} =U_0,\ldots, U_d$ such that 

$$\Bigg\{\begin{array}{l|l} y\in X & \exists \text{ }  f_1, \ldots, f_k\in F\text{, and } x\in U_j\text{ } (0\leq j\leq d) \text{ such that } \\   \text{ } & y = f_k\cdots f_1\cdot x  \text{ and } f_{k_0}\cdots f_1\cdot x\in U_j \text{ } \forall \text{ } 1\leq k_0\leq k \end{array}\Bigg\}$$

is uniformly finite (i.e. the same bound applies independent of $x$).  }\label{dad definition}
\normalfont

The following terminology for DAD makes the analogy with asymptotic dimension clear, and will be critical to the proof of our main theorem.  

\definition{Let $\mathcal{P}_{fs}(\Gamma)$ denote the collection of all finite, symmetric subsets of $\Gamma$ containing the identity.  Let $\Gamma\acts X$ be an action by homeomorphisms, $B\subset X$ a subset, and $F\in \mathcal{P}_{fs}$.  An $F$-chain in $B$ is a finite sequence $x_0, \ldots, x_n$ of points in $B$ such that $x_{i+1} = f_{i+1}\cdot x_i$ with $f_{i+1}\in F$ ($0\leq i<n$).  Two points $x, y\in B$ are in the same $F$-component of $B$ if they are connected by an $F$-chain in $B$ (this is an equivalence relation).  We say $B$ is $F$-connected if any two points in $B$ can be connected by an $F$-chain.  If $S\in \mathcal{P}_{fs}(\Gamma)$, we say $B\subset X$ is $S$-bounded if it is a subset of $S\cdot x$ for some $x\in X$.  A cover of $X$ by $m+1$ (open/closed) sets whose $F$-components are $S$ bounded is called a (open/closed) $(m, F, S)$-cover for $\Gamma\acts X$.  Notice that in this case $S$ can always be taken to be $F^r$ for some $r>0$.  An $n$-dimensional control function for an action $\Gamma\acts X$ on a compact space by homeomorphisms is a function $\mathcal{D}_{X}: \mathcal{P}_{fs}(\Gamma)\to \mathcal{P}_{fs}(\Gamma)$ such that for every $F\in \mathcal{P}_{fs}(\Gamma)$, there is an $(n, F, \mathcal{D}_X(F))$-cover for $\Gamma\acts X$.  Note that $\dad_{free}(\Gamma\acts X)\leq n$ iff there is an $n$-dimensional control function for $\Gamma\acts X$.  } \label{chain definition}

\remark{Although the terminology in \ref{chain definition} relates specifically to the definition of $\dad$ for free actions, we will use it even for actions which are not free.  The Hurewicz-type theorem that we will prove in section 4 will be stated using the definition for free actions.  Note, however, that by the definition for non-free actions, any action with infinite stabilizers will be infinite dimensional.  Since the two definitions are equivalent for an action with uniformly-finite stabilizers, and actions with uniformly-finite stabilizers can easily be related to free actions by ignoring some torsion elements, we do not really lose anything by restricting our attention in this way.  From now on we will only use the definition for free actions and will not mention the other definition again until section 6.  } \label{freeness remark}

\normalfont
For actions on totally disconnected spaces, we can equivalently require the covers to be clopen, as shown below.  

\lemma{Suppose $\mathcal{V} = \{V_0, \ldots, V_d\}$ is an open cover of a totally disconnected, compact space $X$.  Then there is a clopen cover $\mathcal{W} = \{W_0, \ldots, W_d\}$ such that $W_i\subset V_i$ for all $0\leq i\leq d$.  In the context of the previous definition, this implies if $\mathcal{V}$ is a $(d, F, S)$-cover, then $\mathcal{W}$ is a $(d, F, S)$-cover.}   \label{lemma2}

\begin{proof}

Given $\mathcal{V}$, there is a cover $\mathcal{U} = \{U_i\}_{i=0}^d$ by open sets such that $\overline{U_i}\subset V_i$ for $0\leq i\leq d$ (a more general statement is shown in \cite{munkres}, Lemma 41.6).  Then since $X$ has a basis of clopen sets, we can find for each $i$ and each $x\in \overline{U}_i$ a clopen set $W_{i, x}$ containing $x$ and contained in $V_i$.  For each $i$, the $W_{i, x}$'s are a cover of $\overline{U_i}$ (which is compact since $X$ is compact); and so there is a finite subcover, the union over which we denote $W_i$.  This is a finite union of clopen sets and therefore clopen.  Moreover, the collection $\{W_i\}_{i=0}^d$ is a cover of $X$, and $W_i\subset V_i$ for all $i$.  \end{proof} 

\definition{If $\Gamma\acts X$ and $A\subset X$, we write $\Gamma\acts A$ for the partially defined action obtained by restricting $\Gamma\acts X$ to $A$ (more precisely, this is the restriction to $A$ of the transformation groupoid of $\Gamma\acts X$).  The $\dad_{free}$ of this partially-defined action is defined using the same definition as in \ref{chain definition} with covers required to be open in the subspace topology on $A$ and $F$-chains required to be contained in $A$ (this aligns with the definition of $\dad(X\rtimes \Gamma)|_A$ using the groupoid definition \cite[Definition 5.1]{dasdimGWY}).  See \cite[Part I]{Exel2015PartialDS} for a precise definition of partially defined actions.  We will sometimes write $\dad_{free}(A)$ to mean the dimension of the partially-defined action restricted to $A$.  An $n$-dimensional control function for $A\subset X$ is then a function $\mathcal{D}_A: \mathcal{P}_{fs}(\Gamma)\to \mathcal{P}_{fs}(\Gamma)$ such that for every $F\in \mathcal{P}_{fs}(\Gamma)$, there is an open $(n, F, \mathcal{D}_A(F))$-cover of $A$.  }\label{dad of restricted action}

\examples{Two important sources of examples in coarse geometry are graphs, considered as metric spaces by equipping their vertex set with the $l^1$-path metric, and coarse disjoint unions of bounded spaces.  Groups can be considered as metric spaces in the finitely generated case through the Cayley graph construction.  More generally, any countable group has a right-invariant proper metric unique up to coarse equivalence \cite[Lemma 2.3.3]{notes_on_prop_A}.  }

\definition{If $(G_n)$ is a countable collection of finite metric spaces, the coarse disjoint union $\sqcup_{n} G_n$ is the metric space with underlying set the disjoint union of the $G_n$ and metric $d(x, y):=d_{G_n}(x, y)$ if $x, y\in G_n$ and $d(x, y):= \diam(G_n) + \diam(G_m)$ if $x\in G_n$ and $y\in G_m$ with $m\neq n$.  If $\Gamma$ is a group and $(N_i)$ is a collection of finite-index subgroups, we define $\Box_{(N_i)}\Gamma := \sqcup_i \Gamma/N_i$ where each $\Gamma/N_i$ has the metric induced from a fixed, right-invariant proper metric on $\Gamma$ (i.e. the distance between cosets is the infimum over the distances between their different elements according to the original metric on $\Gamma$).  This is called a box space of $\Gamma$.  If the metric on $\Gamma$ is the word metric coming from some generating set and if the $N_i$ are normal, the induced metrics on quotients are the word metrics coming from the image of that generating set in the quotient.  }

\definition{A filtration of $\Gamma$ is a collection $(N_i)$ of finite index subgroups of $\Gamma$ such that $\bigcap_i N_i = \{e\}$, the trivial subgroup.  If $\Gamma$ is residually finite, the collection of all finite index subgroups is then a filtration.  Note that a countable, residually finite group has a filtration which is also a sequence.  }\label{filtration definition}

\remark{Our definition of coarse disjoint union is slightly different from the one found in \cite[Section 2]{boxspacesDT}, and so the box space we associate to $(N_i)$ is not necessarily a coarse disjoint union over $\Gamma/N_i$ in their sense (as the space between different components may not tend to infinity).  However, our definition of box space is the same as theirs in the case when $\Gamma$ is finitely generated and residually finite and $(N_i)$ is a filtration of $\Gamma$.  Note also that we do not always require the subgroups $N_i$ to be normal, as this is not necessary for some of our results.  } %generalize boxspacesDT stuff
%(and appears to be unnecessary for the main results of \cite{boxspacesDT}

\definition{If $(G_i)$ is a sequence of finite metric spaces, we say $\asdim(G_i)\leq d$ uniformly if for all $r>0$ there is a $R>0$ and a $(d, r, R)$-cover of each $G_i$.  We similarly say for a sequence $(\Gamma\acts X_i)$ of actions on finite spaces that $\dad(\Gamma\acts X_i)\leq d$ uniformly if for all $F\in \mathcal{P}_{fs}(\Gamma)$ there is $S\in \mathcal{P}_{fs}(\Gamma)$ and a $(d, F, S)$-cover for each action.  }  

\lemma{$\asdim\Box_{(N_i)}\Gamma\leq d$ if and only if $\asdim(\Gamma/N_i)\leq d$ uniformly.  }\label{uniform and box space dimension}

\begin{proof} For the first inequality, notice that a $(d, r, R)$-cover for $\Box_{(N_i)}\Gamma$ gives rise to a $(d, r, R)$-cover for each $\Gamma/N_i$ (just take the intersection). 

For the second inequality, let $r>0$ and find a $(d, r, R)$-cover $\mathcal{U}^i = \{U^i_0, \ldots, U^i_d\}$ of each $\Gamma/N_i$.  Let $I_r = \{i\mid \diam(\Gamma/N_i)\leq r\}$.  Then $\cup_{i\in I_r} \Gamma/N_i\subset \Box_{(N_i)}\Gamma$ has diameter at most $2r$ (hence is $0$-dimensional) and $\Box_{(N_i)}\Gamma\setminus \big(\cup_{i\in I_r}\Gamma/N_i)$ is $r$-disjoint.  The desired inequality follows from \cite[Theorem 1]{MR1808331}.  \end{proof}

\lemma{Suppose $G = \bigcup_i G_i$ is an increasing union of groups and $X  = \Box_{(H_j)} G = \bigcup_i X_i$ is an increasing union of box spaces where $X_i = \Box_{(H_j\cap G_i)}G_i$ is a box space of $G_i$ (here the metrics on $X$ and $G_i$ are induced from the right invariant, proper metric on $G$, and the metric on $X_i$ is constructed using the metric on $G_i$).  Then $\asdim G = \sup_i \asdim G_i$ and $\asdim X = \sup_i \asdim X_i$.  }\label{direct limit lemma}

\begin{proof}
Fix $R>0$.  Then there is $N$ such that $B_e^R(G)\subset G_N$.  Notice that $d_G(G_N\cdot g, G_N\cdot h)$ is either $0$ or $>R$.  This follows since if $d_G(G_N\cdot g, G_N\cdot h)\leq R$ then the two orbits have a nontrivial intersection and therefore coincide.  Since different $G_N$-orbits are $>R$ apart, we can cover each one separately using $\asdim G_N + 1$-sets.  For the second claim, find $N$ as before.  Then if $x, y\in G/H_j$ are in separate $G_N$-orbits, their representatives in $G$ must be in separate $G_N$ orbits, and so $d_{G/{H_j}}(x, y)>R$.  We can therefore similarly cover each orbit separately using $\asdim \Box_{(H_j\cap G_N)} G_N + 1$ sets.  \end{proof}

\normalfont
When the space being acted on is discrete, there is a very concrete relationship between the dynamic asymptotic dimension of the action and the asymptotic dimension of the orbits.  

\lemma{Suppose $\Gamma$ is a countable group with subgroup $N<\Gamma$ and $|\Gamma|$ is $\Gamma$ equipped with a right-invariant, proper metric.  Suppose further that $|\Gamma/N|$ is a quotient of $\Gamma$ with the metric induced from $|\Gamma|$.  Consider the action $\Gamma\acts \Gamma/N$ on left cosets by left multiplication (this is an action even if $N$ is not normal).  Then an $S$-bounded subset has diameter at most $\diam(S)$ in $|\Gamma|$.  Moreover, a subset $A\subset |\Gamma/N|$ with diameter at most $R$ is $B_e^R(|\Gamma|)$-bounded (recall that quotients do not increase the growth of a group).  In particular, this implies a $(d, B_e^r(|\Gamma|), S)$-cover for $\Gamma\acts \Gamma/N$ is a $(d, r, \diam(S))$-cover of $|\Gamma/N|$; and a $(d, r, R)$-cover of $|\Gamma/N|$ is a $(d, B_e^r(|\Gamma|), B_e^R(|\Gamma|))$-cover for $\Gamma\acts \Gamma/N$.  }\label{basic lemma} 

\begin{proof}
Observe that $x, y\in A\subset \Gamma$ are connected by a $B_e^r(|\Gamma|)$-chain in A iff they are connected by an $r$-chain in $A\subset |\Gamma|$.  This follows since $d_{|\Gamma/N|}(\gamma\cdot x, x)\leq r$ for any $\gamma\in B_e^{r}(|\Gamma|)$ and conversely any elements $\leq r$ away from each other in $|\Gamma/N|$ are connected by some $\gamma\in B_e^r(|\Gamma|)$.  \end{proof}

\normalfont
There is a familiar dynamical object which is reminiscent of a box space called an odometer.  We will see in \ref{odometers and box spaces} that the asymptotic dimension of box spaces is closely related to the dynamic asymptotic dimension of odometers.  

\definition{Let $\Gamma$ be a countable group and $(N_i)$ a countably infinite collection of finite-index subgroups directed by inclusion.  The profinite completion of $\Gamma$ along $(N_i)$ is the inverse limit $\displaystyle{\lim_{\leftarrow i} \Gamma/N_i}$ and is denoted $\widehat{\Gamma}_{(N_i)}$.  If $(N_i)$ is the collection of all such subgroups, we just write $\widehat{\Gamma}$.  If the $N_i$ are normal subgroups, there is a group structure on $\widehat{\Gamma}_{(N_i)}$ and an action $\Gamma\acts \widehat{\Gamma}_{(N_i)}$ given by the quotient map and left multiplication.  Otherwise, there is still an action $\Gamma\acts \widehat{\Gamma}_{(N_i)}$ by left multiplication on left cosets in each coordinate.  The canonical quotient map $p_i: \widehat{\Gamma}_{(N_i)}\to \Gamma/N_i$ is continuous for the profinite topology and equivariant.  }

\section{Odometers and Box Spaces} \label{odometers and box spaces}
\normalfont
As we have seen, asymptotic dimension and DAD are closely related at least when the space being acted on is discrete.  This relationship is encapsulated by \ref{basic lemma} and, together with some topological facts about profinite completions, can be used to show the DAD (according to the definition for free actions) of an odometer formed from a sequence $(N_i)$ of subgroups is the same as the asymptotic dimension of the box space formed from $(N_i)$.  Consequently, we see that questions about asymptotic dimension of box spaces are special cases of questions about $\dad$.  This connection is spiritually similar to how the DAD of an action is bounded above by the asymptotic dimension of its warped cone (this can be shown directly, but see \cite{warpedcones} for definitions and a more complete discussion of warped cones and DAD).  

\remark{It will be important for \ref{prop2} that the collection $(N_i)$ is a sequence, meaning the index set is the natural numbers and $N_{i+1}\subset N_i$.  }

\lemma{Let $\Gamma$ be a countable group and $(N_i)$ a countable collection of finite index subgroups which is directed by inclusion.  A subset $U\subset \widehat{\Gamma}_{(N_i)}$ is clopen iff it is of the form $p_i^{-1}(S)$ where $S\subset \Gamma/N_i$ is any subset. } \label{lemma3}

\begin{proof}

First, sets of the form $p_i^{-1}(S)$ are clopen as $S\subset \Gamma/N_i$ is clopen and $p_i$ is continuous.  

Fix the metric $d((x_i), (y_i)) = \sum_{i=1}^{\infty} \frac{d_{\Gamma/N_i}(x_i, y_i)}{2^i\diam(\Gamma/N_i)}$ on $\widehat{\Gamma}$.  If $U$ is clopen, then $\widehat{\Gamma} = U\cup U^c$ where $U$ and $U^c$ are separate connected components.  In fact, since $\widehat{\Gamma}$ is a compact metric space, there is $\epsilon$ such that the $\epsilon$-neighborhoods of $U$ and $U^c$ do not intersect.  

Now, suppose $U$ is not of the desired form.  Then for all $i$ there is $j\geq i$ and $x_j\in\widehat{\Gamma}$ such that $p_j(x_j)\in p_j(U)$ but $p_{j+1}(x_j)\notin p_{j+1}(U)$.  But then $x_j\notin U$ for all such (inifinitely-many) $j$ and $d(x_j, U)\leq \sum_{n=j+1}^\infty \frac{1}{2^n}\to 0$ as $j\to \infty$\footnote{Technically, the sum is over a rearrangement, but that doesn't matter.}.  Hence, there are points in $U^c$ arbitrarily close to $U$.  By what was said above, this implies $U$ is not clopen. \end{proof}

\proposition{Let $\Gamma$ be a countable group and $(N_i)$ a sequence of finite-index subgroups which is directed by inclusion.  Let $\Gamma\acts \widehat{\Gamma}_{(N_i)}$ be the odometer action induced by left multiplication on each finite quotient.  If $\dad_{free}(\Gamma\acts \widehat{\Gamma}_{(N_i)})\leq d$ then $\asdim(\Box_{(N_i)}\Gamma)\leq d$.  } \label{prop2}

%can use free dad right?

%needs to be directed, but doesn't have to be a sequence!?

\begin{proof}
Fix a proper, left-invariant metric on $\Gamma$.  Let $R\geq 0$ be given and set $F = B_e^R(\Gamma)$.  Let $\mathcal{V} = \{V_0, \ldots, V_d\}$ be an open $(d, F, M_F)$-cover for $\Gamma\acts \widehat{\Gamma}_{(N_i)}$.  Observe that for each $i$, the $F$-components of $p_i(V_j)$ are $R$-separated according to the metric on $\Gamma/N_i$.  

%Find $i_R$ such that $d_{\Box_{(N_i)}\Gamma}(\Gamma/N_{i_1}, \Gamma/N_{i_2})\geq R$ whenever $i_1, i_2\geq i_R$ and $i_1\neq i_2$ (assuming the diameter of $\Gamma/N_i$ tends to infinity).  

%Observe that every quotient $\Gamma/N_i$ in the inverse system $(\Gamma/N_i)$ is contained in a subsystem which is also a sequence.  

Now by \ref{lemma2} we can assume each $V_j$ is clopen and hence by \ref{lemma3} is the pullback of $S\subset \Gamma/N_i$ for some $i$.  Find $i_0$ such that all of the $V_j$ are pullbacks of subsets of $\Gamma/N_{i_0}$.  Then if $i\geq i_0$, $p_i(x)\in p_i(V_j)$, and $f\cdot p_i(x)\in p_i(V_j)$, then $x, f\cdot x\in V_j$.  Thus, the cardinality of each $F$-component of $p_{i}(V_j)$ is bounded by $M_F$ for all $i\geq i_0$; so by \ref{basic lemma} the diameter (as a subset of $\Box_{(N_i)}\Gamma$) of each $R$-component of $p_i(V_j)\subset |\Gamma/N_i|$ is bounded by $RM_F$ for $i\geq i_0$.  This gives a $(d, R, RM_F)$-cover of each $\Gamma/N_i$ for all $i\geq i_0$.  Since $(N_i)$ is a sequence, the set $\{i \mid i < i_0\}$ is finite, and so the diameter of $\Gamma/N_i$ for $i<i_0$ is uniformly bounded.  We have therefore shown $\asdim(\Gamma/N_i)\leq d$ uniformly.  The last part then follows from \ref{uniform and box space dimension}.   \end{proof}

\proposition{Let $\Gamma$ be a countable group and $(N_i)$ a countably infinite collection of finite index subgroups of $\Gamma$ directed by inclusion.  Let $\Gamma\acts \widehat{\Gamma}_{(N_i)}$ be the corresponding odometer action (given by left multiplication as before).  Then $\dad_{free}(\Gamma\acts \widehat{\Gamma}_{(N_i)})\leq \asdim(\Box_{(N_i)}\Gamma)$.  } \label{other box space inequality}
\begin{proof}

Fix a finite subset $F\subset \Gamma$.  

%The coarse structure on $\Box_{(N_i)}\Gamma$ can be metrized by choosing any proper, left-invariant metric on $\Gamma$ and then taking the quotient metric induced on each finite quotient (see \cite[Remark 3.2 (2)]{finnsell2015asymptotic}).  

Find $R>0$ so that $d_\Gamma(\gamma\cdot x, x) < R$ for all $\gamma\in F$ and $x\in \Gamma$, hence $d_{\Gamma/N_i}(\gamma\cdot x, x)< R$ for all $\gamma\in F$ and all $i$.  Suppose $\asdim\Box_{(N_i)}\Gamma \leq n$.  Then $\asdim(\Gamma/N_i)\leq n$ uniformly.  Let $\mathcal{U}^i = \{U^i_0, \ldots, U^i_n\}$ be a $(R, M_R)$-cover of $\Gamma/{N_i}$.  Then by \ref{basic lemma} each $\mathcal{U}^i$ is also a $(d, F, \#B_{e}^{M_R + 1}(\Gamma/N_i))$-cover for $\Gamma\acts \Gamma/{N_i}$ (and since the cardinality of balls in finite quotients of $\Gamma$ is bounded by the cardinality of balls in $\Gamma$, this bound is independent of $i$).  If $K = \ker(\Gamma\to \widehat{\Gamma}_{(N_i)})$, then $\Gamma/K\acts \widehat{\Gamma}_{(N_i)}$ freely and is orbit equivalent to the original $\Gamma$-odometer.  Find $i_F$ so that the quotient $\Gamma/K\to \Gamma/N_{i_F}$ induces a map $C_F(\Gamma/K)\to C_F(\Gamma/N_{i_F})$ which is an isometry on balls of radius $\leq M_R + 2$.  Then pulling back $\mathcal{U}^{i_F}$ by the projection $p_{i_F}: \widehat{\Gamma}_{(N_i)}\to \Gamma/N_{i_F}$ gives a $(d, F, \#B_e^{M_R + 1}(\Gamma/N_{i_F}))$-cover for $\Gamma\acts \widehat{\Gamma}_{(N_i)}$.  \end{proof}

\normalfont
We now state both inequalities in one theorem:

\theorem{Let $\Gamma$ be a countable group and $(N_i)$ a sequence of finite index subgroups which is directed by inclusion.  Suppose $\Gamma\acts \widehat{\Gamma}_{(N_i)}$ by left multiplication.  Then $\dad_{free}(\Gamma\acts \widehat{\Gamma}_{(N_i)}) = \asdim\Box_{(N_i)}\Gamma$.  } \label{box space theorem} \qed


\corollary{Let $\Gamma$ be a countable, residually finite group and $(N_i)$ a sequence which is also a filtration of $\Gamma$ (see \ref{filtration definition}).  Then $\dad(\Gamma\acts \widehat{\Gamma}_{(N_i)}) = \asdim\Box_{(N_i)}\Gamma$.  } \qed

%might need finitely-generated to use D and T stuff

\normalfont
This can also be combined with \cite[3.1]{boxspacesDT}:

\corollary{Let $\Gamma$ be a finitely generated, residually finite group and $(N_i)$ a sequence of normal subgroups which is also a filtration.  Then either $\dad(\Gamma\acts \widehat{\Gamma}_{(N_i)}) = \infty$ or $\dad(\Gamma\acts \widehat{\Gamma}_{(N_i)}) = \asdim \Gamma$.  The dimension is finite if $\Gamma$ has polynomial growth (equivalently is virtually nilpotent).  } \qed

\normalfont
Finally, we can use the relationship between the dimensions odometers and box spaces to obtain an extension theorem for the asymptotic dimension of box spaces which has been desired for several years.  

\corollary{Suppose $1\to \Delta\to \Gamma\xrightarrow{\pi} \Lambda\to 1$ is an exact sequence of countable groups.  Suppose that $(N_i)$ is a sequence of finite index subgroups of $\Gamma$ directed by inclusion.  This sequence induces sequences and hence box spaces of $\Delta$ and $\Lambda$.   Then if $\asdim\Box_{(N_i\cap \Delta)}\Delta = m$ and $\asdim\Box_{(\pi(N_i))}\Lambda = n$, $\asdim\Box_{(N_i)}\Gamma \leq m + n$.  } \label{box spaces theorem} 

\section{Product Theorem}
\normalfont
Now we examine permanence properties of $\dad$.  Most basically, one would expect DAD to be subadditive over products of groups acting on products of spaces, and the proof of this fact will exhibit some of the same ideas used in the next, more technical, section.  We rely heavily on the ideas of \cite{Brodskiy2006AHT} for both this section and the next and have noted when one of our lemmas is inspired by one of theirs.  Indeed, if one relaxes the definition of $\dad$ to allow Borel sets, or even sets obtained from finite operations with open sets, the work of \cite{Brodskiy2006AHT} can be transported to the dynamical setting with very little extra work.  Our solution to the extra technicality created by the topology of $X$ comes from a lemma established in this section which we refer to as simply `the trick'.  This lemma capitalizes on some rigidity to provide a way of switching between covers using open or closed sets while preserving the properties related to $\dad$.  

\definition{Let $\Gamma\acts X$ be an action on a compact space by homeomorphisms and $A\subseteq X$.  If $k\geq n+1\geq 1$, an $(n, k)$-dimensional control function for $A$ is a function $\mathcal{D}_{A}:\mathcal{P}_{fs}(\Gamma)\to \mathcal{P}_{fs}(\Gamma)$ such that for any $F\in \mathcal{P}_{fs}(\Gamma)$ there is a $(k-1, F, \mathcal{D}_{A}(F))$-cover $\mathcal{U} = \{U_0, \ldots, U_{k-1}\}$ of $A$ by sets open in the subspace topology on $A$ such that each $x\in A$ belongs to at least $k-n$ elements of $\mathcal{U}$ (equivalently $\cup_{i\in T} U_i = A$ for every $T\subset \{0, \ldots, k-1\}$ consisting of $n+1$ elements). } \label{control function space definition}

\lemma{(the trick) Let $\Gamma\acts X$ be an action on a compact space by homeomorphisms and $A\subset X$ closed.  
\begin{itemize}
\item[(1)] Suppose $F, S\in\mathcal{P}_{fs}(\Gamma)$, $n>0$, $k\geq n+1$, and $\mathcal{U} = \{U_0, \ldots, U_{k-1}\}$ is a $(k-1, F, S)$-cover for $A$ by open sets such that each $x$ belongs to $\geq k-n$ elements.  Then there is a $(k-1, F, S)$-cover by sets which are the closures of open sets and still have the property that each $x$ belongs to $\geq k-n$ elements.  

\item[(2)] If $\Gamma\acts X$ is either free or isometric, and if $U$ is a closed set whose $F$-components are $S$-bounded, then $U$ is contained in an open set whose $F$-compoenents are $S$-bounded.  
\end{itemize} } \label{control functions lemma}

\begin{proof}
For the first part, let $\mathcal{U} = \{U_0, \ldots, U_{k-1}\}$ be such a cover.  Then $\bigcup_{S\subset \{0, \ldots, k\} : |S| = k-n} \cap_{s\in S}U_s$ is an open cover, and so has some Lebesgue number $\lambda$.  This shows every $x\in X$ has a $\lambda$-ball which is contained in at least $k-n$ elements of $\mathcal{U}$.  Replace the $U_i$ with the closures of their $\lambda/2$-interiors.  

For the second part, let $U$ be a closed set.  For $x\in U$, denote by $Y_x$ the set $\{y \mid y \text{ is connected to $x$ by an } F^r\text{-chain in $U$}\}$.  Then there is $\epsilon>0$ such that 

$$\inf_{x\in U} \inf_{Y_x} \inf_{\{\gamma\in F \mid \gamma\cdot y\notin U\}} d(\gamma\cdot y, U) >\epsilon$$
where we interpret the innermost infimum as $\infty$ when $\{\gamma\in F^r \mid \gamma\cdot y\notin U_i\}$ is empty.  This follows from a compactness argument and the properties of $\mathcal{U}$.  

Now assume $\Gamma\acts X$ is free.  Let $S\subset \Gamma$ be the union over all $x\in U$ of the sets $S_x = \{\gamma\in \Gamma \mid \gamma\cdot x = y\in Y_x\}$.  Since the $Y_x$ are uniformly finite and $\Gamma\acts X$ is free, $S$ is finite.  Find $0<\delta<\epsilon/2$ such that $d(\gamma\cdot x, \gamma\cdot y)<\epsilon/2$ whenever $d(x, y)<\delta$ for all $\gamma\in FS$.  We can then replace $U$ with its open $\delta$-neighborhood.   

If instead $\Gamma\acts X$ is isometric, we can simply take $\delta = \epsilon/2$.  \end{proof}

\remark{We will frequently apply the lemma above by simply saying `use the trick', when we want to replace an open cover by a closed one or a closed set with a larger open set.  }

%\lemma{Suppose $\Gamma\acts X$ and $A\subset X$ is such that for each $r>0$ there is a $}

\lemma{(compare to \cite[Theorem 2.4]{Brodskiy2006AHT}) Suppose $\Gamma\acts X$ is either isometric or free and $A\subset X$ is closed.  If $\mathcal{D}_{A}^{(n+1)}$ is an $n$-dimensional control function for $A$, $F\in \mathcal{P}_{fs}(\Gamma)$, and we define $\{\mathcal{D}_{A}^{(i)} \}_{i\geq n+1}$ inductively on powers of $F$ by $\mathcal{D}_{A}^{(i+1)}(F^r) = B_e^r(C_F(\Gamma))\cdot \mathcal{D}^{(i)}_{A}(F^{3r})$; then for each $k$, there is an open $(k-1, F^r, \mathcal{D}^{(k)}_{A}(F^r))$-cover of $A$ such that each $x\in A$ is covered by at least $k - n$ elements of the cover (we say it this way since the functions $\mathcal{D}_{A}^{(k)}$ are not defined on all elements of $\mathcal{P}_{fs}(\Gamma)$ if $\Gamma$ is not finitely generated, so we can't quite call them control functions).  }\label{extra cover lemma}

\begin{proof}

The case $k=n+1$ is obvious.  Suppose the result holds for some larger $k$.  Let $\mathcal{U} = \{U_1, \ldots, U_k\}$ be a open $(k-1, F^{3r}, \mathcal{D}_{A}^{(k)}(F^{3r}))$-cover for $\Gamma\acts X$ with all $x\in A$ in at least $k-n$ elements of $\mathcal{U}$.  Using the trick (1), we can replace the sets $U_i$ by the closures of open sets $V_i$ which form a cover of $A$ with the same properties.  Define $U_i' = (F^r\cdot U_i)\cap A$ (these are closed sets).  Then $\mathcal{U}' = \{U_1', \ldots, U_k'\}$ is a closed $(k-1, F^r, B_e^r(C_F(\Gamma))\cdot \mathcal{D}_{A}^{(k)}(F^{3r}))$-cover for $A$.  Define $U_{k+1}'$ to be the union of all the (disjoint, open) sets $W_S := (\cap_{s\in S} V_s)\setminus \cup_{i\notin S} U_i'$ where $S\subset \{1, \ldots, k\}$ has $k-n$ elements.  We claim that for $S\neq T$ (where $S, T\subset \{1, \ldots, k\}$ each have $k-n$ elements) that the sets $W_S$ and $W_T$ are $F^r$-disjoint (i.e. not connected by an $F^r$-chain contained in their union).  

Suppose $a, b\in U_{k+1}'$ are such that $a\in \cap_{t\in T} V_t \setminus \cup_{i\notin T} U_i'$ and $b\in \cap_{s\in S} V_s \setminus \cup_{i\notin S} U_i'$ for $T\neq S$ and $|S| = |T| = k-n$.  Then there is $t\in T\setminus S$ such that $a\in V_t$, so if $f\cdot a = b$ for some $f\in F^r$, then $b\in U_t'\supseteq F^r\cdot V_t$, a contradiction.  

Now suppose $x\in A$ belongs to exactly $k-n$ sets $U_i'$, $i\leq k$.  Since the sets $\cap_{s\in S}V_s$ form a cover as $S$ ranges over subsets of $\{1, \ldots, k\}$ with $|S| = k-n$; $x\in \cap_{s\in S}V_s$ for some $s\in S$.  But then if $x\notin U_{k+1}'$, we must have $x\in U'_j$ for some $j\notin S$, contradicting that $x$ is only contained in $k-n$ sets $U_i'$.  

Thus, $\{U'_{k+1}, U'_k, \ldots, U'_1\}$ is a $(k, F^r, B_e^r(C_F(\Gamma))\cdot \mathcal{D}_{A}^{(k)}(3r))$-cover with each $x\in A$ contained in at least $k-n + 1$ sets (recall that $U_{k+1}'$ is open and the other sets are closed).  We can then apply the trick (2) to each closed set to produce a cover by open sets with the same properties.  \end{proof} 

%prove product theorem here

\theorem{(compare to \cite[Theorem 2.5]{Brodskiy2006AHT}) If $\Gamma\acts X$ and $\Lambda\acts Y$ are free, then $\text{DAD}(\Gamma\times \Lambda\acts X\times Y)\leq \text{DAD}(\Gamma\acts X) + \text{DAD}(\Lambda\acts Y)$.  

\begin{proof}

Fix $F\in\mathcal{P}_{fs}(\Gamma)$ and $H\in\mathcal{P}_{fs}(\Lambda)$.  Suppose $\Gamma\acts X$ and $\Lambda\acts Y$ have dimension at most $m$ and at most $n$, respectively, and let $\mathcal{D}_X$ and $\mathcal{D}_Y$ be their respective control functions.  Let $k = m + n + 1$.  Use \ref{extra cover lemma} to produce the functions $\mathcal{D}^{(k)}_{X}$ and $\mathcal{D}^{(k)}_{Y}$.  So there is an open $(k-1, F, \mathcal{D}^{(k)}_{\Gamma\acts X}(F))$-cover $\mathcal{U} = \{U_i\}_{i=0}^{k-1}$ of $X$ and an open $(k-1, H, \mathcal{D}^{(k)}_{\Lambda\acts Y}(H))$-cover $\mathcal{V} = \{V_i\}_{i=0}^{k-1}$ of $Y$ such that each $x\in X$ is in at least $k-m$ elements of $\mathcal{U}$ and each $y\in Y$ is in at least $k-n$ elements of $\mathcal{V}$.  Since $k-m + k - n = k+1$, the family $\{U_i\times V_i\}_{i=0}^{k-1}$ is a cover of $X\times Y$.  Moreover, it is a $(F\times H, \mathcal{D}^{(k)}_{\Gamma\acts X}(F) \cdot \mathcal{D}^{(k)}_{\Lambda\acts Y}(H))$-cover for $\Gamma\times \Lambda\acts X\times Y$.  \end{proof}

\section{Hurewicz-type Theorem}
\normalfont
In this section, we will prove a more general permanence property witnessing a natural dimension-theoretic property of $\dad$.  We start by defining the dimension of an appropriate morphism in the category of group actions.  As dynamic asymptotic dimension is meant to measure the large scale complexity of orbits, the appropriate kind of morphism here is one which sends orbits to orbits and induces a coarse map of the orbit structures.  The following definition is not new, but it was inspired by the definition of continuous orbit couple in \cite[Definition 2.6]{LiDynamics}.  A continuous orbit couple induces a `coarse equivalence of orbits', so we weaken this definition to one which will give a `coarse map' of orbits.  

\definition{Suppose $\Gamma\acts X$ and $\Lambda\acts Y$ are actions on compact spaces by homeomorphisms and that $f: X\to Y$ and $\alpha: \Gamma \times X\to \Lambda$ are continuous maps with $f(\gamma\cdot x) = \alpha(\gamma, x)\cdot x$.  Define $\dad(f) = \sup_{A\subset Y : \dad(A) = 0} \dad_{free}(f^{-1}(A))$ with $\dad(f^{-1}(A))$ defined as in \ref{dad of restricted action}.  }\label{dad of a map}

\normalfont
The intuition behind this definition is we want to measure the largest dimension of the partially defined action on a subspace of $X$ which $f$ collapses down to something zero dimensional in the image.  Recall for instance that in the category of vector spaces, this would be the nullity of a linear map, i.e. the dimension of the pullback of $\{0\}$.  In the category of topological spaces, one can similarly define the dimension of a function to be the supremum over the dimensions of pullbacks of points.  However, as in the case of asymptotic dimension, we must alter this definition slightly by considering pullbacks of arbitrary $0$-dimensional subspaces.  See  \cite[Proposition 4.1]{Brodskiy2006AHT} for an example which shows this deviation is necessary.  

%\remark{The $\dad$ of a restricted (partially defined action) is the same as the usual definition with covers required to be open in the subspace topology.  }

%\definition{An $m$-dimensional control function for $\Gamma\acts X$ is a function $\mathcal{D}_X: \mathcal{P}_{fs}(\Gamma)\to \mathcal{P}_{fs}(\Gamma)$ such that for all $F\in \mathcal{P}_{fs}(\Gamma)$, there is an open $(m, F, \mathcal{D}_X(F))$-cover for $\Gamma\acts X$.  Note that $\dad(\Gamma\acts X)\leq d$ iff there exists a $d$-dimensional control function for $\Gamma\acts X$.  }

\normalfont
Many of the ideas in this section come from \cite{Brodskiy2006AHT}, so as in the previous section we note at the beginning of each statement what lemma is being analogized.  

\definition{Let $\Gamma\acts X$, $\Lambda\acts Y$ and $f: X\to Y$ be as in \ref{dad of a map}.  An $m$-dimensional control function for $f$ is a function $\mathcal{D}_f: \mathcal{P}_{fs}(\Gamma)\times \mathcal{P}_{fs}(\Lambda)\to \mathcal{P}_{fs}(\Gamma)$ such that for all $S\in \mathcal{P}_{fs}(\Gamma)$ and $T\in \mathcal{P}_{fs}(\Gamma)$, we have that if $B\subset Y$ is $T$-bounded, then $f^{-1}(B)$ has a $(m, S, \mathcal{D}_f(S, T))$-cover (open in the subspace topology on $f^{-1}(B)$).  }

\lemma{(see \cite[Proposition 4.5]{Brodskiy2006AHT}) Let $\Gamma\acts X$, $\Lambda\acts Y$ and $f: X\to Y$ be as in \ref{dad of a map}.  If $\dad(f)\leq m$, then there is an $m$-dimensional control function for $f$.  }\label{dimension implies function}

\begin{proof}
Fix $T\in \mathcal{P}_{fs}(\Lambda)$ and $S\in \mathcal{P}_{fs}(\Gamma)$.  Suppose for a contradiction that $F_n$ is an increasing sequence of sets in $\mathcal{P}_{fs}(\Gamma)$ with $\cup_n F_n = \Gamma$, and that for each $n$, $B_n = T\cdot x_n$ is a $T$-bounded set such that $f^{-1}(B_n)$ does not have an open $(m, S, F_n)$-cover.  Then for any subsequence $n_k$, $f^{-1}(\bigcup_{k} B_{n_k})$ must have $\dad$ greater than $m$.  

Notice that at least one of the following alternatives must hold: either there exist infinitely many $n$ with $x_n$ in different orbits, or there exist infinitely many $n$ with $x_n$ in the same orbit.  In the first case, we get a subsequence $n_k$ with each $x_{n_k}$ in a different orbit so that $\bigcup_{k} B_{n_k}$ is zero-dimensional; this is a contradiction.  In the second case, we have a subsequence $n_k$ with each $x_{n_k}$ in the same orbit, say the orbit of $x_0$.  Fix a right-invariant, proper metric $d_\Lambda$ on $\Lambda$.  Then define a proper metric $d$ on the orbit of $x_0$ by the formula $d(x, y):=\inf_{\{\lambda\in \Lambda : \lambda\cdot x = y\}} d_\Lambda(\lambda, e)$.  Then either $d(x_{n_k}, x_0)\to \infty$ or these distances do not tend to infinity.  In the case of the latter, that implies $\bigcup_{k} B_{n_k}$ is finite and hence zero dimensional, a contradiction.  In the case of the former, passing to a further subsequence yields a sequence $n_l$ with $d(x_{n_l}, \bigcup_{j<l} x_{n_j})\to \infty$.  Again, this implies $\bigcup_{l} B_{n_l}$ is zero dimensional, which is a contradiction.  \end{proof}

%\definition{An $(m, k)$-dimensional control function for $f:X\to Y$ is a function $\mathcal{D}_f: \mathcal{P}_{fs}(\Gamma)\times \mathcal{P}_{fs}(\Lambda)\to \mathcal{P}_{fs}(\Gamma)$ such that for all $S\in \mathcal{P}_{fs}(\Gamma)$ and $T\in \mathcal{P}_{fs}(\Gamma)$, we have that if $B\subset Y$ is $T$-bounded, then $f^{-1}(B)$ has a $(k-1, S, \mathcal{D}_f(S, T))$-cover and every $x\in f^{-1}(B)$ is covered by at least $k-m$ sets.}  

\lemma{Suppose $X$ and $Y$ are compact metric spaces and $f: X\to Y$ is continuous.  Let $C\subset Y$ be closed and $N$ an open neighborhood of $f^{-1}(C)$.  Then there is an open neighborhood $G$ of $C$ such that $f^{-1}(G)\subset N$.  } \label{neighborhoods lemma}

\begin{proof}
Using the metric on $Y$, one can construct a family $\{G_\alpha\}$ of open sets containing $C$ such that $\bigcap_\alpha \overline{G}_\alpha = C$ ($\overline{G}_\alpha$ denotes the closure of $G_\alpha$).  Then $\bigcap_\alpha f^{-1}(\overline{G}_\alpha) = f^{-1}(\bigcap_\alpha \overline{G}_\alpha) = f^{-1}(C)$.  That is, for all $x\in f^{-1}(C)^c$, there is $\alpha$ such that $x\notin f^{-1}(\overline{G}_\alpha)$.  Moreover, since $f^{-1}(\overline{G}_\alpha)$ is closed, $d(x, G_{\alpha})>\delta_x$ for some $\delta_x>0$ 

Now let $N\supset f^{-1}(C)$ be an open neighborhood.  Then $N^c$ is closed and therefore compact.  For each $x\in N^c\subset f^{-1}(C)^c$, let $\delta_x>0$ as above.  Then the collection of open balls $\{B_x^{\delta_x}(X)\}$ is a cover of $N^c$ and therefore has a finite subcover $\{B_{x_k}^{\delta_{x_k}}(X)\}_{k=1}^K$.  Let $G_{\alpha_k}$ be the neighborhoods from the paragraph above coming from $x_k$.  Then if we put $G = \bigcap_k G_{\alpha_k}$, $G$ is an open neighborhood of $C$ and $f^{-1}(\bigcap_{\alpha_k}G_{\alpha_k}) = \bigcap_{\alpha_k} f^{-1}(G_{\alpha_k})\subset N$.  \end{proof}

%\lemma{Suppose $X$ and $Y$ are compact, totally disconnected metric spaces and that $f: X\to Y$ is continuous.  Let $C\subset Y$ be closed and $N$ an open neighborhood of $f^{-1}(C)$.  Then there is a neighborhood $G$ of $C$ such that $f^{-1}(G)\subset N$.  } \label{neighborhoods lemma}
%
%\begin{proof}
%Let $\{G_\alpha\}$ be the collection of all clopen neighborhoods of $C$.  Since $Y$ is normal and has a basis of clopen sets, $\bigcap_\alpha G_\alpha = C$ and so $\bigcap_\alpha f^{-1}(G_\alpha) = f^{-1}(\bigcap_\alpha G_\alpha) = f^{-1}(C)$.  That is, for all $x\in f^{-1}(C)^c$, there is $\alpha$ such that $x\notin f^{-1}(G_\alpha)$.  Moreover, since $f^{-1}(G_\alpha)$ is closed, there is $\delta_x>0$ such that $d(x, f^{-1}(G_\alpha))<\delta_x$.  
%
%Now let $N\supset f^{-1}(C)$ be an open neighborhood.  Then $N^c$ is closed and therefore compact.  For each $x\in N^c\subset f^{-1}(C)^c$, let $\delta_x>0$ as in the above.  Then the collection of open balls $\{B_x^{\delta_x}(X)\}$ is a cover of $N^c$ and therefore has a finite subcover $\{B_{x_k}^{\delta_{x_k}}(X)\}_{k=1}^K$.  Let $G_{\alpha_k}$ be the neighborhoods from the paragraph above coming from $x_k$.  Then if we put $G = \bigcap_k G_{\alpha_k}$, $G$ is an open neighborhood of $C$ and $f^{-1}(\bigcap_{\alpha_k}G_{\alpha_k}) = \bigcap_{\alpha_k} f^{-1}(G_{\alpha_k})\subset N$.  \end{proof}

\remark{To prove the main result, it is enough for \ref{neighborhoods lemma} to hold for $C$ a finite subset.  }

\lemma{(compare to \cite[Proposition 4.7]{Brodskiy2006AHT}) Let $\Gamma\acts X$, $\Lambda\acts Y$, and $f: X\to Y$ be as in \ref{dad of a map} and that each action is either free or isometric.  Suppose $\mathcal{D}_f: \mathcal{P}_{fs}(\Gamma)\times \mathcal{P}_{fs}(\Lambda)\to \mathcal{P}_{fs}(\Gamma)$ is an $m$-dimensional control function for $f$.  If we define $\{\mathcal{D}_f^{(i)}\}_{i\geq m+1}$ inductively on powers of $F\in \mathcal{P}_{fs}(\Gamma)$ by $\mathcal{D}_f^{(i+1)}(F^r, T):= B^r_e(C_F(\Gamma))\cdot \mathcal{D}_f^{(i)}(F^{3r}, T)$, then for each $k$ and $T$-bounded subset $B\subset Y$, there is an open $(k-1, F^r, \mathcal{D}_f^{(k)}(F^r, T))$-cover of $f^{-1}(B)$ such that each $x\in X$ is covered by at least $k-m$ elements of the cover.  }\label{function extra cover lemma}

\begin{proof}
The function $F\mapsto \mathcal{D}_f(F, T)$ is an $m$-dimensional control function for $f^{-1}(B)$; more precisely, it is a control function for \textit{any} $T$-bounded subset $B$.  We can therefore apply \ref{extra cover lemma} to produce $(m, k)$-dimensional control functions for $f^{-1}(B)$ (technically these functions are only defined on powers of $F$, but that's all we need to establish this lemma).  Since the construction in the proof of \ref{extra cover lemma} depends only on the control function we start with, these new control functions are the same for any $T$-bounded subset $B\subset Y$, and so we have the desired conclusion.  \end{proof}


\lemma{(see \cite[Proposition 4.8]{Brodskiy2006AHT}) Let $\Gamma\acts X$, $\Lambda\acts Y$ and $f: X\to Y$ be as in \ref{dad of a map} and also assume each action is either free or isometric.  Suppose $\mathcal{D}_f$ is an $m$-dimensional control function for $f$.  Then for all $T, S\in \mathcal{P}_{fs}(\Lambda)$, $F\in \mathcal{P}_{fs}$, and $B\subset Y$ an open set with $T$-components $S$-bounded, $f^{-1}(B)$ can be covered by $k$ sets which are open in $X$ and whose $(F, T)$-components are $\mathcal{D}^{(k)}_{f}(F, S)$-bounded (with $\mathcal{D}_f^{(k)}$ defined as in \ref{function extra cover lemma}), and every $x\in f^{-1}(B)$ is contained in at least $k-m$ elements of this cover.  } \label{control function tricky lemma}

\begin{proof}
Let $C_{\alpha}$ be a $T$-component of $B$.  By assumption, $C_{\alpha} = S_\alpha\cdot x_\alpha$ where $S_\alpha\subset S$, so by \ref{function extra cover lemma}, $f^{-1}(C_\alpha)$ has a $(k-1, F, \mathcal{D}^{(k)}_f(F, S))$-cover by sets which are open in the subspace topology on $f^{-1}(C_\alpha)$ and which has the additional property that every element of $f^{-1}(C_\alpha)$ is contained in at least $k-m$ elements of the cover.  Since $f^{-1}(C_\alpha)$ is closed and hence compact, we can use the trick (1) to replace this cover with a cover by closed sets (which are closed in $X$), and then use the trick (2) to again replace the cover with a cover by sets which are open in $X$, all the while preserving the fact that every element of $f^{-1}(C_\alpha)$ is contained in at least $k-m$ elements of this new cover.  Call this cover $\mathcal{U}^\alpha = U^\alpha_0, \ldots, U^\alpha_{k-1}$.  It covers not just $f^{-1}(C_\alpha)$, but an open neighborhood of it which we can assume without loss of generality to be contained in $f^{-1}(B)$.  By \ref{neighborhoods lemma}, it is possible to find an open neighborhood $G_\alpha$ of $C_\alpha$ that $f^{-1}(G_\alpha)$ is covered by $\mathcal{U}^\alpha$.  

%Since $B$ is closed, we can use the trick (2) to find a neighborhood $N_{\alpha}$ of $C_{\alpha}$ (not necessarily contained in $B$) having the form $\sqcup_{s\in S_\alpha} s\cdot N_{\epsilon_\alpha}(x_0)$ for some open $\epsilon_\alpha$-neighborhood of $x_0$; and such that the $T$-components of $N_{\alpha}$ are $S_\alpha$-bounded).  
%; and such that if $x, y\in N_{\alpha}$ with $\gamma\cdot x = y$ and $y$ can be reached from $x$ by a $T$-chain in $B$, then $\gamma\in S_\alpha$


%Any neighborhood of $C_\alpha$ contained in $B$ will still have $F$-components $S_\alpha$-bounded.  

Find an open neighborhood $M_\alpha$ of $x_\alpha$ such that $N_\alpha := S_\alpha\cdot M_\alpha\subset G_\alpha\subset B$.  Then a $T$-chain beginning in $N_\alpha$ stays in $N_\alpha$ until leaving $B$ (at which point it also won't be in $N_{\alpha'}$ for any $\alpha'\neq \alpha$).  

The collection $\{N_\alpha\}$ covers $B$.  Now let $W_i^{\alpha} = U_i^{\alpha}\cap f^{-1}(N_{\alpha})$ and put $W_i = \bigcup_\alpha W_i^{\alpha}$.  Then $\{W_i\}_{i=0}^{k-1}$ is a cover of $f^{-1}(B)$ by $k$ open sets.  Since any $x\in f^{-1}(B)$ is in $f^{-1}(C_\alpha)$ for some $\alpha$, every $x\in f^{-1}(B)$ is contained in at least $k-m$ elements of $\{W_i\}_{i=0}^{k-1}$.  Moreover, since $f(W_i^{\alpha})\subset N_{\alpha}$, the properties of $N_\alpha$ imply that an $(F, T)$-component of $W_i$ is contained in an $F$-component of some $W_i^{\alpha}\subset U_i^{\alpha_j}$ and is hence $\mathcal{D}^{(k)}_f(F, S)$-bounded. \end{proof}

%We want to make this cover disjoint in a particular way: first remove from each $D_{\alpha_j}$ $(j>1)$ any overlap with $D_{\alpha_1}$, then do the same for $D_{\alpha_2}$ and so on.  This lets us replace $\{D_{\alpha_j}\}_{j=1}^J$ with a cover $\{E_j\}_{j=1}^J$ of $B$ by disjoint clopen sets such that $E_j\subset D_{\alpha_j}$; and such that if $x\in E_j$, $y\in E_{j'}$ with $j\neq j'$, then $x$ and $y$ are not connected by a $T$-chain in $B$ (the cover $\{D_{\alpha_j}\}_{j=1}^J$ had this property, and our way of disjointifying preserves it).  


\definition{Let $f: X\to Y$ be as in \ref{dad of a map}.  Suppose that $F\in \mathcal{P}_{fs}(\Gamma)$ and $S\in \mathcal{P}_{fs}(\Lambda)$.   An $(F, S)$-chain in $U$ is an $F$-chain whose image under $f$ is an $S$-chain.  Two points $x_1, x_2\in U$ are in the same $(F, S)$-component if they are connected by an $(F, S)$-chain.  We say a subset $A\subset X$ is $(F, S)$-bounded if $A$ is $F$-bounded and $f(A)$ is $S$-bounded.  } \label{more chain definition}

\lemma{(Finite union lemma, compare to \cite[Corollary 3.6]{Brodskiy2006AHT}) With the same setup as \ref{dad of a map}, let $A, B\subset X$, $F\in \mathcal{P}_{fs}(\Gamma)$ and $S\in \mathcal{P}_{fs}(\Lambda)$.  Assume that the $(F^{r_X^A}, S^{r^A_Y})$-components of $A$ are $(F^{R_X^A}, S^{R_Y^A})$-bounded, that the $(F^{r_X^B}, S^{r^B_Y})$-components of $B$ are $(F^{R_X^B}, S^{R_Y^B})$-bounded; and that $R_X^B + 2r_X^B<r_X^A$ and $R_Y^B + 2r_Y^B< r_Y^A$.  Then the $(F^{r_X^B}, S^{r^B_Y})$-components of $A\cup B$ are $$(B_e^{2r_X^A + R_X^A}(C_F(\Gamma)), B_e^{2r_Y^A + R_Y^A}(C_S(\Lambda)))\text{-bounded}$$

Here, $C_{F}(\Gamma)$ denotes the Cayley graph of $\Gamma$ with respect to the set $F$ (which need not generate all of $\Gamma$).  }\label{finite union lemma}
\begin{proof}

Let $x_0, \ldots, x_n$ be an $(F^{r_X^B}, S^{r_Y^B})$-chain in $A\cup B$ with no repeated points.  Suppose $x_j$ and $x_k$ are two consecutive points contained only in $A$.  Then $x_{j+1}, \ldots, x_{k-1}$ is an $(F^{r_X^B}, S^{r_Y^B})$-chain in $B$ and therefore is $(F^{R_X^B}, S^{R_Y^B})$-bounded.  

We therefore have that $x_{j+1}$ and $x_{k-1}$ are in the same $(F^{R_X^B}, S^{R_Y^B})$-component, so $x_i$ and $x_j$ are in the same $(F^{R_X^B + 2r_X^B}, S^{R_Y^B + 2r_Y^B})$-component of $A$.  Since $r_X^BR_X^B + 2r_X^B < r_X^A$ and $r_Y^BR_Y^B + 2r_Y^B < r_Y^A$, the points in the original chain which are contained only in $A$ therefore form an $(F^{r_X^A}, S^{r_Y^A})$-chain in $A$ and so (considered as a subset of $A$) are $(F^{R_X^A}, S^{R_Y^A})$-bounded.  The original chain, considered as a set, is therefore $(F^{2R_X^B + 2r_X^B + R_X^A}, S^{2R_Y^B + 2r_Y^B + R_Y^A})$-bounded.  This implies the $(F^{r_X^B}, S^{r_Y^B})$-components of $A\cup B$ are $(B_e^{2R_X^B + 2r_X^B + R_X^A}(C_F(\Gamma)), B_e^{2R_Y^B + 2r_Y^B + R_Y^A}(C_S(\Lambda)))$-bounded, which implies the result since $2R_X^B + 2r_X^B< 2r_X^A$ and $2R_Y^B + 2r_Y^B<2r_Y^A$.  \end{proof}

\theorem{(compare to \cite[Theorem 4.9]{Brodskiy2006AHT}) Let $k = m + n + 1$, where $m, n\geq 0$.  Let $\Gamma\acts X$, $\Lambda\acts Y$ and $f: X\to Y$ be as in \ref{dad of a map} with each action either free or isometric.  Suppose $\mathcal{D}_{Y}$ is an $n$-dimensional control function for $\Lambda\acts Y$ and $\mathcal{D}_f$ is an $m$-dimensional control function for $f$.  Then there is an $m+n$-dimensional control function for $\Gamma\acts X$ depending only on $\mathcal{D}_{Y}$, $\mathcal{D}_f$, $m$, and $n$.  } \label{control function theorem}

\begin{proof}Let $\alpha: \Gamma\times X\to \Lambda$ be the function described in \ref{dad of a map} for $f$.  Apply \ref{extra cover lemma} to $D_{Y}$ to produce the function $\mathcal{D}^{(k)}_{Y}$.   We can assume without loss of generality that $\mathcal{D}^{(k)}_{Y}(T)\supset T$ for all $T\in \mathcal{P}_{fs}(\Lambda)$ (one can always make the outputs of a control function larger, as long as they are still finite subsets).  




%We will construct a finite subset and show that $X$ can be covered by $k$ open sets $\{D^j\}_{j=1}^k$ with all $F$-components of $D^j$ being.  

Fix $F\in \mathcal{P}_{fs}(\Gamma)$.  Define inductively a sequence of finite subsets $S_Y^{(n+1)} \subset T_Y^{(n+1)}\subset S_Y^{(n)}\subset T_Y^{(n)}\subset \cdots \subset S_Y^{(1)}\subset T_Y^{(1)}\subset S_Y^{(0)}\subset T_Y^{(0)}$ starting from $S_Y^{(n+1)} = \alpha(F, X)$ and moving to lower indices so that for every $i$ we have $T_Y^{(i)} = \mathcal{D}^{(k)}_Y(S_Y^{(i)})$ and $S_Y^{(i)} = (T_Y^{(i+1)})^3$.  

Write $Y$ as the union of $n+1$ open sets $\{A_i\}_{i=1}^{n+1}$ such that all $S_Y^{(0)}$-components of $A_i$ are $T_Y^{(0)}$-bounded for every $i$.  Since $A_i\subset Y$ we can use \ref{extra cover lemma} to obtain a cover of $A_i$ by $k$ sets $\{U_i^j\}_{j=1}^k$ which are open in $X$ such that all $S_Y^{(i)}$-components of $U_i^j$ are $T_Y^{(i)}$-bounded for every $j$ and every point $y\in A_i$ belongs to at least $k-n = m+1$ sets.  To be clear, although $A_i$ is not necessarily closed, we can apply \ref{extra cover lemma} to $Y$ and then take intersections with $A_i$.  

Apply \ref{function extra cover lemma} to $\mathcal{D}_f$ and $F$ to obtain the function $\mathcal{D}^{(k)}_f$.     As with $\mathcal{D}_{\Lambda\acts Y}$, we can assume $\mathcal{D}_f(S, T)\supset S$ for all $S\in \mathcal{P}_{fs}(\Gamma)$ and $T\in \mathcal{P}_{fs}(\Lambda)$.  

Define inductively another sequence of finite sets $S_X^{(n+1)} \subset T_X^{(n+1)}\subset S_X^{(n)}\subset T_X^{(n)}\subset \cdots S_X^{(1)}\subset T_X^{(1)}\subset S_X^{(0)}\subset T_X^{(0)}$ starting with $S_X^{(n+1)} = F$ and for every $i$ defining $T_X^{(i)} = \mathcal{D}^{(k)}(S_X^{(i)}, T_Y^{(i)})$ and $S_X^{(i)} = (T_X^{(i+1)})^3$.  

By \ref{function extra cover lemma} and the properties of $A_i$, $f^{-1}(A_i)$ has a cover by $k$ sets $\{B_i^j\}_{j=1}^k$ which are open in $X$ such that all $(S_X^{(i)}, S_Y^{(i)})$-components are $(T_X^{(i)}, T_Y^{(0)})$-bounded for every $j$ and every point $x\in f^{-1}(A_i)$ belongs to at least $k-m = n + 1$ sets.  

Put $D_i^j = B_i^j\cap f^{-1}(U_i^j)$ and let $D^j$ be the union of all the $D_i^j$.  The $D^j$'s cover $X$ by Kolmogorov's argument: given $x\in X$, the set of $j$'s such that $x\in B_i^j$ has at least $n + 1$ elements, and the set of $j$'s such that $f(x)\in U_i^j$ has at least $k-n$ elements, so these sets of such $j$'s cannot be disjoint.  

Notice all $(S_X^{(i)}, S_Y^{(i)})$-components of the set $D_i^j$ are $(T_X^{(i)}, T_Y^{(i)})$-bounded.  By \ref{finite union lemma}, all $(S_X^{(n+1)}, S_Y^{(n+1)})$-components of the set $D^j$ are $((T_X^{(1)})^3, (T_Y^{(1)})^3)$-bounded.  Since $(S_X^{(n+1)}, S_Y^{(n+1)}) = (F, \alpha(F, X))$ by definition, and $(F, \alpha(F, X))$-components coincide with $F$-components, all $F$-components of $D^j$ are $(T_X^{(1)})^3$-bounded.  \end{proof}

\theorem{(Hurewicz-type theorem) Suppose $\Gamma\acts X$ and $\Lambda\acts Y$ are actions on compact spaces by homeomorphisms which are each either isometric or free, and that $f: X\to Y$ and $\alpha: \Gamma\times X\to \Lambda$ are continuous functions such that $f(\gamma\cdot x) = \alpha(\gamma, x)\cdot f(x)$.  With $\dad(f)$ defined as in \ref{dad of a map}, it holds that $\dad_{free}(\Gamma\acts X)\leq \dad(f) + \dad_{free}(\Lambda\acts Y)$.  }\label{hurewicz-type theorem} \qed

\remark{The product theorem proven in the previous section can now be obtained as a consequence of \ref{hurewicz-type theorem}.  In fact, one can show a bit more by considering actions on bundles which respect the bundle structure.  The map from the bundle to its base space will have dynamic asymptotic dimension exactly that of the action on a fibre.  }

\corollary{Suppose $\Gamma\acts X$, $\Lambda\acts Y$, and $f: X\to Y$ are as above.  Then $f$ has an $m$-dimensional control function if and only if $\dad(f)\leq m$.  }\label{function iff dimension}
\begin{proof}
One direction was already proved in \ref{dimension implies function}.  For the other direction, apply \ref{control function theorem} to each subspace of $Y$ such that the restricted action is $0$-dimensional. \end{proof}

\section{Applications}

%\lemma{Suppose $\Gamma$ is residually finite and $(N_i)$ is a collection of finite index subgroups which is a filtration of $\Gamma$ and is directed by inclusion.  If $1\to \Delta\to \Gamma\to \Lambda\to 1$ is an extension of groups, $(N_i)$ induces collections of finite index subgroups $(N_i\cap \Delta)$ of $\Delta$ and $(N_i/(\Delta\cap N_i))$ of $\Lambda$.  If $\dad(\Delta\acts \widehat{\Delta}_{(N_i\cap \Delta)})\leq m$, and $f: \widehat{\Gamma}_{(N_i)}\to \widehat{\Lambda}_{(N_i/(\Delta\cap N_i)}$ is the map induced by the quotients $\Gamma/N_i\to \Lambda/(N_i/\Delta\cap N_i)$, then there is an $m$-dimensional control function for $f$.  }

\normalfont
%Our first application shows that $\dad$ of odometers behaves subadditively over extensions of groups.  This gives a tool for producing many new examples of actions with finite $\dad$, as any group built from extensions of groups with finite dimensional odometers will itself have a finite dimensional odometer.  

An important corollary of most Hurewicz-type theorems is that dimension is subadditive over extensions of objects.  Although there are natural statements of this kind that we can now prove for the dynamic asymptotic dimension, it is less obvious whether such theorems can be usefully applied.  For instance, it is possible to construct actions by $\Z^2$ where the restriction to either copy of $\Z$ is minimal so that the action of the other copy on the quotient by the orbit of the first is infinite dimensional (because the quotient is very non-Hausdorff).  Notwithstanding, we highlight two useful applications below.  The first is that the dynamic asymptotic dimension does not change when passing to a finite index subgroup, which shows an additional aspect in which $\dad$ behaves like a large-scale dimension for group actions.  The second is an extension theorem for the dimension of odometer actions, which will have further applications in the final two sections.  

For some applications, it is easier to directly produce a control function rather than bound the dimension of a map.  In fact, this inspires an alternative definition for said dimension.  

\definition{Suppose $\Gamma\acts X$, $\Lambda\acts Y$, and $f: X\to Y$ are as in \ref{dad of a map}.  The value $\dad(f)$ can equivalently be defined to be the smallest $m$ such that $\dad(f^{-1}(\{y\}))\leq m$ uniformly for all $y\in Y$.  }
\begin{proof}
If $\dad(f)\leq m$ according to the original definition, then $f$ has an $m$-dimensional control function by \ref{function iff dimension}.  This implies $\dad(f^{-1}(\{y\}))\leq m$ uniformly.  

Now assume $\dad(f^{-1}(\{y\}))\leq m$ uniformly.  Then using \ref{finite union lemma} and induction produces an $m$-dimensional control function for $f$.  \end{proof}

\corollary{Suppose $\Gamma\acts X$ is a free action on a compact metric space by homeomorphisms and that $\Lambda<\Gamma$ is a finite index subgroup.  Then $\dad(\Gamma\acts X) = \dad(\Lambda\acts X)$.  }\label{new dad of a map}
\begin{proof}
By passing to the normal core, we may assume without loss of generality that $\Lambda$ is normal.  Suppose $\dad(\Lambda\acts X) = m$.  Let $X/\Lambda$ be the quotient of $X$ by the orbit equivalence relation coming from $\Lambda\acts X$ and let $\Gamma/\Lambda\acts X/\Lambda$ be the obvious action.  Let $f: X\to X/\Lambda$ be the quotient map.  Since $\Gamma/\Lambda\acts X/\Lambda$ is locally finite, $\dad(\Gamma/\Lambda\acts X/\Lambda) = 0$.  Moreover, if $y\in X/\Lambda$, $f^{-1}(\{y\})$ is a $\Lambda$-orbit in $X$ and so, as $\Gamma\acts X$ is free, $\dad(\Gamma\acts f^{-1}(\{y\}))\leq \dad(\Lambda\acts f^{-1}(\{y\}))\leq \dad(\Lambda\acts X)$ uniformly for all $y$, and so $\dad(\Gamma\acts f^{-1}(\{y\}))\leq m$ uniformly.  \end{proof}

\normalfont
%For our second application, we need one more lemma connecting some ideas from Section 3.  This will allow us to avoid having to assume the sequence $(N_i)$ is normal.  
%
%\lemma{Suppose $\Gamma\acts \widehat{\Gamma}_{(N_i)}$ is an odometer and $\Delta<\Gamma$ is a subgroup.  Then $\dad(\Delta\acts \widehat{\Gamma}_{(N_i)})\leq d$ if and only if $(\Delta\acts \Gamma/N_i)\leq d$ uniformly.  }\label{uniform lemma}
%
%\begin{proof}The `if' direction follows from \ref{basic lemma} and \ref{other box space inequality}.  The `only if' direction follows from \ref{prop2} and \ref{basic lemma}.   \end{proof}
% need some more details ^^^^^^

%\corollary{Suppose $\Gamma\acts X$ and $\Lambda\acts Y$ are free actions on compact metric spaces by homeomorphisms.  Suppose $E = \bigsqcup_{x\in X} \{(y, x) \mid y\in Y\}$ is a fibre bundle with base $X$ and fibre $Y$.  There is a free action $\Gamma\times \Lambda\acts E$ given by  }

\corollary{Suppose $\Gamma$ is residually finite and $(N_i)$ is a collection of finite index normal subgroups which is a filtration of $\Gamma$ and is directed by inclusion.  If $1\to \Delta\to \Gamma\to \Lambda\to 1$ is an extension of groups, $(N_i)$ induces collections of finite index subgroups $(N_i\cap \Delta)$ of $\Delta$ and $(N_i/(\Delta\cap N_i))$ of $\Lambda$.  Then $\dad_{free}(\Gamma\acts \widehat{\Gamma}_{(N_i)})\leq \dad_{free}(\Delta\acts \widehat{\Delta}_{(N_i\cap \Delta)}) + \dad_{free}(\Lambda\acts \widehat{\Lambda}_{(N_i/(\Delta\cap N_i))})$.  }

\begin{proof}
Supposing that $\dad(\Delta\acts \widehat{\Delta}_{(N_i\cap \Delta)})\leq m$ and $f: \widehat{\Gamma}_{N_i}\to \widehat{\Lambda}_{(N_i/(\Delta\cap N_i))}$ is obtained from the quotient maps $\Gamma/N_i\to \Lambda/(N_i/(\Delta\cap N_i))$, we will show $\dad(f)\leq m$ using \ref{new dad of a map}.  Note that the action $\Lambda\acts \widehat{\Lambda}_{(N_i/(\Delta\cap N_i))}$ need not be free, but is isometric.  

Let $y\in \widehat{\Lambda}_{(N_i/(\Delta\cap N_i))}$.  Then $f^{-1}(\{y\})$ is the $\Delta$-orbit of some $x\in \widehat{\Gamma}_{(N_i)}$.  Right multiplication by $x^{-1}$ (using that the $N_i$ are normal) shows the action $\Delta\acts f^{-1}(\{y\})$ is continuously orbit equivalent to the action of $\Delta$ on $\mathcal{O}_\Delta((eN_i)_i)$, that is, the $\Delta$-orbit of the identity element in $\widehat{\Gamma}_{N_i}$.  Moreover, the latter action is continuously orbit equivalent to the action of $\Delta$ on the orbit of $(eN_i)_i$ in $\widehat{\Delta}_{(N_i\cap \Delta)}$.  

Since $(N_i)$ is a filtration of $\Gamma$, any $\gamma\notin \Delta$ moves a point in the $\Delta$-orbit of $x$ outside of this $\Delta$-orbit.  Then if $\Gamma\acts f^{-1}(\{y\})$ is the restricted action, we have $\dad(\Gamma\acts f^{-1}(\{y\})) = \dad(\Delta\acts f^{-1}(\{y\}) = \dad(\Delta\acts \mathcal{O}_{\Delta}((eN_i)_i))\leq \dad(\Delta\acts \widehat{\Delta}_{(N_i\cap \Delta)})$.  In fact, these equalities and inequalities hold quantitatively and independently of $y$, so we have $\dad(f)\leq m$.   \end{proof}


%Let $B\subset \widehat{\Lambda}_{(N_i/(\Delta\cap N_i))}$ have the form $S\cdot y$ for some $S\in \mathcal{P}_{fs}(\Lambda)$.  Fix $F\in \mathcal{P}_{fs}(\Gamma)$, and let $y\in \widehat{\Lambda}_{(N_i/(\Delta\cap N_i))}$.  Notice that $\Delta\acts f^{-1}(\{y\})$ is isomorphic to $\Delta\acts\widehat{\Delta}_{(N_i\cap \Delta)})$.  Moreover, as $(N_i)$ is a filtration of $\Gamma$, $\gamma\cdot x\notin f^{-1}(\{y\})$ for any $x\in f^{-1}(\{y\})$ and $\gamma\in \Gamma\setminus \Delta$.  This means that a cover for $\Delta\acts f^{-1}(\{y\})$ whose $F$-components are $T$-bounded will still have this property as a cover for $\Gamma\acts f^{-1}(\{y\})$.  We can therefore use that $\dad(\Delta\acts \widehat{\Delta}_{(N_i\cap \Delta)})\leq m$, together with the finite union lemma \ref{finite union lemma} to show a cover for $f^{-1}(B)$ whose $F$-components are bounded with a bound depending on $F$ and $S$ (more precisely the cardinality of $S$).

%Let $f: \widehat{\Gamma}_{(N_i)}\to \widehat{\Lambda}_{(N_i/(\Delta\cap N_i))}$ be the continuous map induced by the quotient maps $\Gamma/N_i\to \Lambda/(N_i/(\Delta\cap N_i))$.  If $y\in \widehat{\Lambda}_{(N_i/(\Delta\cap N_i))}$, $f^{-1}(\{y\})$ is a $\Delta$-orbit.  Since $\Gamma$ is residually finite, $\gamma\cdot x\notin f^{-1}(\{y\})$ for any $x\in f^{-1}(\{y\})$ and $\gamma\notin \Delta$.  Thus, $\dad(\Gamma\acts f^{-1}(\{y\}))\leq \dad(\Delta\acts f^{-1}(\{y\}))\leq \dad(\Delta\acts \widehat{\Gamma}_{(N_i)})$.  
%
%If we have a $(d, F, S)$-cover for $\Delta\acts \Delta/N_i$, we can produce a $(d, F, S)$-cover for $\Delta\acts \Gamma/N_i$: just use the same cover for each $\Delta$ orbit (which is isomorphic to $\Delta\acts \Delta/N_i$), which works because the different orbits don't interact.  If $\dad_{free}(\Delta\acts \widehat{\Delta}_{(N_i\cap \Delta)}) \leq d$, then $\dad_{free}(\Delta\acts \Delta/(N_i\cap \Delta))\leq d$ uniformly.  But then this implies $\dad(\Delta\acts \Gamma/(N_i\cap \Delta))\leq d$ uniformly, and so $\dad(\Delta\acts \widehat{\Gamma}_{(N_i)})\leq d$ by \ref{uniform lemma}.  Thus $\dad(\Gamma\acts f^{-1}(\{y\}))\leq \dad(\Delta\acts \widehat{\Delta}_{(N_i\cap \Delta)})$.  Since the cover we get by applying \ref{uniform lemma} does not depend on $y$, we in fact have $\dad(\Gamma\acts f^{-1}(\{y\}))\leq \dad(\Delta\acts \widehat{\Delta}_{(N_i\cap \Delta)})$ uniformly for all $y$.  




%Using the finite union lemma \ref{finite union lemma} as in the previous proof shows there is an $m$-dimensional control function for $f$, and so $\dad(\Gamma\acts X) = m$.  

%can make quotient isometric?

%\theorem{Suppose $\Gamma\acts X$ is a free action on a compact space.  Let $U\subset X$ be an open, invariant subset.  Then $\dad(\Gamma\acts X)\leq \dad(\Gamma\acts U) + \dad(\Gamma\acts U^c)$.  }



\section{Box Spaces of Elementary Amenable Groups}
\normalfont Our extension theorem for the asymptotic dimension of box spaces can be used to bound the dimension of box spaces of many elementary amenable groups in terms of their Hirsch length.  To do this, we need a certain perspective on elementary amenable groups to generalize the Hirsch length, and prove our final results.  

\definition{Define the class of (countable) \textit{elementary amenable} groups to be the smallest class of groups containing all finite and all countable abelian groups which is closed under taking subgroups, quotients, extensions, and direct limits.  }

\begin{lemma}\cite{MR1143191}*{page 238} If $\mathcal{C}$ and $\mathcal{D}$ are classes of groups, denote by $\mathcal{C}\mathcal{D}$ the groups which fit into an extension of the form $1\to C\to \cdot \to D\to 1$ with $C\in \mathcal{C}$ and $D\in \mathcal{D}$ and denote by $L\mathcal{C}$ the groups which are locally in $\mathcal{C}$ (i.e. those groups whose finitely generated subgroups are all in $\mathcal{C}$).  Now define classes of groups $\Gamma_\alpha$ as follows.  Begin with $\Gamma_0 = \{1\}$ and $\Gamma_1$ the class of finitely generated virtually abelian groups.  Assuming the definition of $\Gamma_\alpha$ for $\alpha$ a countable ordinal, define $\Gamma_{\alpha+1}$ to be $L(\Gamma_\alpha)\Gamma_1$.  If $\alpha$ is a countable limit ordinal, then $\Gamma_\alpha:=\bigcup_{\beta<\alpha}\Gamma_\beta$.  The class of countable elementary amenable groups is then $\bigcup_\alpha \Gamma_\alpha$ where $\alpha$ ranges over all countable ordinals.  \label{elementary amenable lemma} \end{lemma} 

\definition{The Hirsch length of a finitely generated group with finite index abelian subgroup $A$ is defined to be $\text{rank}(A)$ \cite{hillman_1991}.  As in \cite{hillman_1991}, we can use \ref{elementary amenable lemma}, to extend the definition of Hirsch length inductively by defining, for a sequence $1\to \Lambda\to \Gamma\to \Delta\to 1$ that $h(\Gamma) = h(\Lambda) + h(\Delta)$; and defining, for an increasing union of groups $\Gamma = \bigcup_n \Gamma_n$ that $h(\Gamma) = \sup_n h(\Gamma_n)$.  It is shown in \cite{hillman_1991} that this is well defined and that Hirsch length satisfies the following
\begin{description}
\item[(1)] for $\Lambda\leq \Gamma$, $h(\Lambda)\leq h(\Gamma)$
\item[(2)] for $\Lambda \mathrel{\unlhd} \Gamma$, $h(\Gamma) = h(\Lambda) + h(\Gamma/\Lambda)$.  
\item[(3)] $h(\Gamma) = \sup\{h(\Lambda) | \Lambda<\Gamma \text{ is finitely generated }\}$  
\end{description}}

\normalfont We can now use the extension theorem proved in the previous section shows that box spaces (and odometers) of many elementary amenable groups is bounded above by the Hirsch length of the group.  Of course, this means any elementary amenable group built up from extensions by groups of finite Hirsch length with have finite-dimensional box spaces.  Such groups include groups which are not virtually nilpotent, as we will show with examples in the last section.   Of course, the theorems given below can be rephrased in terms of the dynamic asymptotic dimension of odometer actions.  The main ideas for the following applications come from the end of \cite[Section 4]{finnsell2015asymptotic}.  

\theorem{Let $\Gamma$ be a countable, residually finite, elementary amenable group.  Then for any sequence  $(N_i)$ of finite index normal subgroups which is also a filtration, $\asdim \Gamma\leq \asdim\Box_{(N_i)} \Gamma\leq h(\Gamma)$.  } \label{hirsch length upper bound}
\begin{proof}
Since $\Gamma$ is residually finite by \cite[Theorem C]{MR285589}, $\asdim\Gamma\leq \asdim \Box_{(N_i)} \Gamma$ by \cite[Proposition 3.1]{boxspacesDT}.  

Using the description of elementary amenable groups given in \ref{elementary amenable lemma}, we can assume $\Gamma$ is in the class $\Gamma_\alpha$ for some countable ordinal $\alpha$.  The case where $\alpha=1$ follows from \cite[Theorem 3.5]{MR2213160} and \cite[Corollary 4.2]{boxspacesDT} (for virtually abelian groups, the dimension of the full box space coincides with the dimension of the group and with its Hirsch length).  Now suppose $\alpha>1$ and assume the second inequality holds for groups in the class $\Gamma_{\alpha-1}$.  Then \ref{direct limit lemma} shows the result holds for groups in $L(\Gamma_{\alpha-1})$; and \ref{box spaces theorem} shows the result holds for groups in the class $\Gamma_\alpha$.  If $\beta$ is a limit ordinal, $\Gamma_\beta = \cup_{\alpha<\beta} \Gamma_\alpha$, and we can apply \ref{direct limit lemma} again.  \end{proof}

\theorem{If $\Gamma$ is countable and residually finite and $\Gamma = \Lambda\rtimes \Delta$ where $\Lambda$ is virtually polycyclic and $\Delta$ is locally finite, then for any sequence of finite index normal subgroups $(N_i)$ which is also a filtration, $\asdim\Gamma = \asdim\Box_{(N_i)}\Gamma = h(\Gamma)$.  }

\begin{proof}
That $\asdim\Gamma \leq h(\Gamma)$ comes from the previous theorem.  The reverse inequality comes from \cite[3.5]{MR2213160} and \cite[58]{asdimBD}, as these show $h(\Gamma) = h(\Lambda) = \asdim\Lambda\leq \asdim\Gamma$ \end{proof}

%wreath products corollary

\corollary{Let $F$ be a finite abelian group and $H$ a virtually polycyclic group.  Denote by $F\wr H$ the reduced wreath product.  Then for any sequence of finite index normal subgroups $(N_i)$ which is also a filtration, $\asdim(F\wr H) = \asdim(\Box_{(N_i)}(F\wr H)) = h(F\wr H) = h(H).$ }
\begin{proof}
The wreath product $F\wr H$ is residually finite by \cite[3.2]{MR87652}.  Furthermore, \cite[Theorem 1]{hillman_1991} shows that $h(F\wr H) = h(H)$.  The corollary then follows by applying the previous theorem.  \end{proof}

\section{Baumslag-Solitar groups}
\normalfont
We conclude by demonstrating how the theoretical machinery introduced in this paper can be used to calculate the dimension of box spaces for many new examples.  We will show that for all integers $n$, the box spaces (coming from sequential filtrations) of the Baumslag-Solitar groups $\text{BS}(1, n)$ have dimension $2$.  This class includes infinitely-many groups which are not virtually nilpotent, providing concrete evidence for the conjecture that finite-dimensional amenable groups should have finite-dimensional box spaces.    

\definition{The Baumslag-Solitar group $\text{BS}(m, n)$ is the group given by the presentation $\langle a, b \mid ba^mb^{-1} = a^n\rangle$.  }

\normalfont
We will calculate the asymptotic dimension of $\text{BS}(1, n)$ for all $n$ and use that to calculate the dimension of its box space.  To do this, we will use \cite{MR2431020}, which shows a finitely presented group of asymptotic dimension $1$ is virtually free by utilizing a relationship between asymptotic dimension and ends of groups.  This theorem has also been proven independently by a number of authors.  See, for example, \cite[Theorem 66]{asdimBD} and the surrounding citations.  We use it here to give an affirmative answer to \cite[Question 1.4]{asdimnote}.  

\corollary{A finitely-presented, torsion free group of asymptotic dimension $1$ is free.  }\label{asdim free}
\begin{proof}
Combine \cite[Theorem 2]{MR2431020} with \cite[Theorem 3]{MR223439}, the latter of which shows that a finitely generated, virtually free, torsion free group is free.  \end{proof}
%
%
\example{The group $\text{BS}(1, n)$ is not virtually nilpotent for $n\geq3$ and has $\asdim\Box_{(N_i)} \text{BS}(1, n) = \asdim \text{BS}(1, n) = 2$ for any sequence $(N_i)$ of finite index normal subgroups which is also a filtration.  } \label{example}

\begin{proof}
The group $\text{BS}(1, n)$ is residually finite by \cite[Theorem C]{MR285589} and fits into the exact sequence $1\to \Z[1/n]\to \text{BS}(1, n)\to \Z\to 1$ (the first map takes $1/n$ to $a$, and the second is the quotient by the subgroup generated by $a$).  So \cite[Theorem 68]{asdimBD} (subadditivity over extensions) shows that $\asdim\text{BS}(1, n) \leq 2$.  Because $\text{BS}(1, n)$ is torsion-free but not free, it cannot be that $\asdim\text{BS}(1, n) = 1$.  

Next, since $\Z[1/n]\cong \Z$, $\Box_{(\Z[1/n]\cap N_i)}\Z[1/n]$  has asymptotic dimension $1$.  By \ref{box spaces theorem}, $\asdim \Box_{(N_i)}  \text{BS}(1, n)\leq 2$.  Since $\text{BS}(1, n)$ is residually finite, $\asdim BS(1, n)\leq \asdim \Box \text{BS}(1, n)$, and so $\asdim \Box_{(N_i)} \text{BS}(1, n) = 2$.  It is shown in \cite{exponentialgrowthBSgroups} that, for $n\geq 3$, $\text{BS}(1, n)$ has exponential growth and is therefore not virtually nilpotent.  \end{proof}

\normalfont
We can generalize the above examples by defining more complicated groups with a similar presentation.  
A further application of the extension theorem gives an upper bound on the dimension of box spaces for these groups.  

\example{Let $m_1, m_2\geq 1$ be integers, $n_{i, j}\in \Z$ for all $1\leq i\leq m_1$ and $1\leq j\leq m_2$, and $\sigma_j\in S_{m_1}$ for $1\leq j\leq m_2$ (here $S_{m_1}$ denotes permutations of $\{1, \ldots, m_1\}$).  

Now define $$\Gamma := \langle a_1, \ldots, a_{m_1}, b_1, \ldots, b_{m_2} \mid [a_i, a_j] = 1, \text{ } [b_i, b_j] = 1, \text{ } b_ja_ib_j^{-1} = a_{\sigma_j(i)}^{n_{i, j}}\rangle.$$

\noindent Then $\asdim \Box_{(N_i)} \Gamma\leq m_1 + m_2$ for any sequence $(N_i)$ of finite index normal subgroups which is also a filtration.  }

\begin{proof}


The group $\Gamma$ is residually finite by \cite[Theorem 1]{MR110750} and fits into the exact sequence $1\to \Z^{m_1}\to \Gamma\to \Z^{m_2}\to 1$.  To be precise, the first map identifies $\Z^{m_1}$ with the subgroup generated by the $a_i$'s, and this subgroup is moreover normal in $\Gamma$.  Therefore, $\asdim\Box_{(N_i)} \Gamma\leq m_1 + m_2$.  \end{proof}

\normalfont There are also groups known as generalized Baumslag-Solitar groups \cite{DELGADO2011398}.  Most of these are non-amenable, and so will have infinite dimensional box spaces, but considering these groups leads to at least one more interesting example.  

\example{Let $\Gamma = \langle a, b \mid a^2 = b^{2}\rangle$.  Then for any sequence $(N_i)$ of finite index normal subgroups which is also a filtration, $\asdim \Box_{(N_i)}\Gamma = 2$.  }

\begin{proof}
Since $\Gamma/\langle a^{2n}\rangle\cong \Z/2n*\Z/2n$ is residually finite, it follows that $\Gamma$ is residually finite. For the upper bound, observe that $\Gamma$ fits into the exact sequence $1\to \Z\to \Gamma\to \Z/2*\Z/2\to 1$ where the first map is the inclusion of the normal subgroup $\langle a^2\rangle$ and the second map is the quotient by this subgroup.  Since $\Z/2*\Z/2$ is the infinite dihedral group (which is virtually cyclic), $\asdim \Box_{(N_i)}\Gamma \leq 2$.  For any element $\gamma\in \Gamma$, there is $n$ such that the class of $\gamma$ in $\Gamma/\langle a^{2n}\rangle$ is non-trivial.   We are therefore done if we show $\asdim \Gamma = 2$.  Since $\Gamma$ is torsion free but not free, this follows from \ref{asdim free}.  \end{proof}

%\example{For $m, k\in \Z$, the group $\text{BS}(m, km)$ has $\asdim\Box_{(N_i)}\text{BS}(m, km) = \asdim\text{BS}(m, km) = 2$ for any sequence $(N_i)$ of finite index subgroups which is also a filtration.  For $k\geq 3$, $\text{BS}(m, km)$ is not virtually nilpotent.  }
%
%\begin{proof}
%Notice that $\langle a^m\rangle$ is a normal subgroup of $\text{BS}(m, km)$ and that $\text{BS}(m, km)/\langle a^m\rangle$ is virtually cyclic, as $\langle b\rangle$ is a finite index subgroup.  We have the exact sequence $1\to \langle a^m\rangle\to \text{BS}(m, km)\to \text{BS}(m, km)/\langle a^m\rangle\to 1$, so the extension theorems for asymptotic dimension of groups and box spaces show $\asdim\text{BS}(m, km)\leq 2$ and $\asdim \Box_{(N_i)}\text{BS}(m, km)\leq 2$.  For the reverse inequality, observe first that $\text{BS}(m, km)/\langle a^{nm}\rangle$ is virtually cyclic for all $n\geq 1$.  This shows $\text{BS}(m, km)$ is residually finite, as any nontrivial element is detected by this quotient for some $k$, and virtually cyclic groups are residually finite.  Since $\text{BS}(m, km)$ is torsion free, but not free, so it cannot be that $\asdim\text{BS}(m, km) = 1$.  
%
%For $k\geq 3$, \cite{exponentialgrowthBSgroups} shows that $\text{BS}(1, k)$ has exponential growth and is therefore not virtually nilpotent.  Since $\text{BS}(1, k)$ is a subgroup of $\text{BS}(m, km)$, the latter is also not virtually nilpotent for $k\geq 3$.  \end{proof}


%\example{The group $\text{BS}(-m, m)$ has $\asdim\Box_{(N_i)}\text{BS}(-m, m) = \asdim\text{BS}(-m, m) = 2$ for any sequence $(N_i)$ of finite index subgroups which is also a filtration.  }
%
%\begin{proof}
%Notice that $\langle a^m\rangle$ is a normal subgroup of $\text{BS}(-m, m)$ and that $\text{BS}(-m, m)/\langle a^m\rangle$ is virtually cyclic, as $\langle b\rangle$ is a finite index subgroup.  We have the exact sequence $1\to \langle a^m\rangle\to \text{BS}(-m, m)\to \text{BS}(-m, m)/\langle a^m\rangle\to 1$, so the extension theorems for asymptotic dimension of groups and box spaces shows $\asdim\text{BS}(-m, m)\leq 2$ and $\asdim \Box_{(N_i)}\text{BS}(-m, m)\leq 2$.  For the reverse inequality, observe first that $\text{BS}(-m, m)/\langle a^{km}\rangle$ is virtually cyclic for all integers $k$.  This shows $\text{BS}(-m, m)$ is residually finite, as any nontrivial element is detected by this quotient for some $k$, and virtually cyclic groups are residually finite.  Like before, $\text{BS}(-m, m)$ is torsion free, but not free, so it cannot be that $\asdim\text{BS}(-m, m) = 1$.  \end{proof}









%Define inductively a sequence $r_Y^{(n+1)} < R_Y^{(n+1)} < r_Y^{(n)} < R_Y^{(n)} < \cdots < r_Y^{(1)} < R_Y^{(1)} < r_Y^{(0)} < R_Y^{(0)}$ of numbers beginning with $r_Y^{(n+1)} = 1$, $R_Y^{(n+1)} = \mathcal{D}^{(k)}_{\Lambda\acts Y}(h(F))$, and then $r_Y^{(i)} = R_Y^{(i+1)} + 2$ ($0\leq i\leq n$), and $R_Y^{(i)} = (R_Y^{(i+1)} - 1)(\mathcal{D}^{(k)}_{\Lambda\acts Y}(h(F)^{r_y^{(i)}}) + 1) + 1$.  
%
%
%%Define inductively a sequence $r_Y^{(n+1)} < R_Y^{(n+1)} < r_Y^{(n)} < R_Y^{(n)} < \cdots < r_Y^{(1)} < R_Y^{(1)} < r_Y^{(0)} < R_Y^{(0)}$ of numbers beginning with $r_Y^{n+1} = 1$, $R_Y^{i} = D_{\Lambda\acts Y}(F^{r_Y^{(i)}})$ and $r_Y^{(i)} = 3R_Y^{(i+1)}$.  
%
%Express $Y$ as the union of $n+1$ sets $\{A_i\}_{i = 1}^{n+1}$ such that all $h(F)^{r_Y^{(0)}}$-components of each $A_i$ have cardinality at most $R_Y^{(0)}$ for every $i$ (this is possible since $R_Y^{(0)}$ is at least $\mathcal{D}_{\Lambda\acts Y}(h(r_Y)^{(0)})$).  Since $A_i\subset Y$ (and by definition of the $r_Y^{(i)}$ and $R_Y^{(i)}$), we can further express $A_i$ as the union of $k$ sets $\{U_i^j\}_{j=1}^{k}$ such that all $h(F)^{r_Y^{(i)}}$-components of each $U_i^j$ have cardinality at most $R_Y^{(i)}$ and every point $y\in A_i$ belongs to at least $k-n = m+1$ sets.  
%
%Apply \ref{inductive prop} to $\mathcal{D}_f$ to produce the function $\mathcal{D}^{(k)}_f$.  Define inductively a sequence of numbers $r^{(n+1)}_X < R^{(n+1)}_X < r^{(n)}_X < R_X^{(n)} < \cdots < r_X^{(1)} < R_X^{(1)}$ starting with $r_X^{(n+1)} = 1$, $R_X^{(n+1)} = \mathcal{D}_f^{(k)}(F, f(F)^{R_Y^{(n+1)}})$ and then $r_X^{(i)} = R_X^{(i+1)} + 2$ and $R_X^{(i)} = (R_X^{(i+1)} - 1)(\mathcal{D}^{(k)}_f(F^{r_X^{(i)}}, f(F)^{r_y^{(i)}R_Y^{(i)}}) + 1) + 1$.  %<<<<<<check this
%
%For every $i$, we can use the properties of $A_i$ and \ref{last lemma} to express $f^{-1}(A_i)$ as the union of $k$ sets $\{B_i^j\}_{j=1}^k$ such that all $(F^{r_X^{(i)}}, h(F)^{r_Y^{(0)}})$-components of $B_i^j$ have cardinality at most $R_X^{(i)}$ for every $j$ and every point $x\in f^{-1}(A_i)$ belongs to at least $k - m = n+1$ sets.  
%
%Put $D_i^j = B_i^j\cap f^{-1}(U_i^j)$ and let $D^j$ be the union of all $D_i^j$ for $i\in \{1, \ldots, n+1\}$.  The collection of the $D^j$ is a cover by Kolmogorov's argument: given $x\in X$ there is $i$ so that $f(x)\in A_i$.  The set of $j$'s such that $x\in B_i^j$ has at least $k-m$-elements and the set of $j$'s such that $f(x)\in U_i^j$ has at least $k-n = m + 1$ elements, so they cannot be disjoint since $k - m + (k-n) = n+1 + m + 1 = k + 1 > k$ and there are $k$ possible $j$'s.    
%
%Notice all $(F^{r_X^{(i)}}, f(F)^{r_Y^{(i)}})$-components of the sets $D_i^j$ have cardinality at most $(R_X^{(i)}, R_Y^{(i)})$.  By repeated application of \ref{finite union lemma} with $r_X^B = r_Y^B = 1$ at every step and $r_X^{A} = r_X^{(n-k)}$ and $r_Y^A = r_Y^{(n-k)}$ at the $k$-th step ($k = 0, \ldots, n$), that is, for the union of $D_0^j$ through $D_k^j$; we can show all $(F^{r_X^{(n+1)}}, f(F)^{r_Y^{(n+1)}})$-components (that is, the $(F, f(F))$-components) of the set $D^j$ have cardinality at most $(R_X^{(1)}, R_Y^{(1)})$.  By definition of $f$, $F$-components and $(F, f(F))$-components coincide, so the $F$-components of $D^j$ have cardinality at most $R_X^{(1)}$ (hence uniformly finite).  Notice that $R_X^{(1)}$ depends only on the control functions $\mathcal{D}_{\Lambda\acts Y}$ and $\mathcal{D}_f$.  









%
%
%
%
%Call this cover $\mathcal{U}^\alpha$.  The neighborhoods we took when using the trick the second time can be as small as we want, so we can ensure the images of these open sets under $f$ will be contained in $N_\alpha$.   The image under $f$ of the union of the $U_i^\alpha$ is therefore a clopen neighborhood of $C_\alpha$ contained in $N_\alpha$, since $f$ is assumed to send clopen sets to clopen sets.  Call this clopen neighborhood $G_\alpha$.  
%
%Repeating this process for every $T$-component of $B$ gives a collection of clopen subsets, $\{G_\alpha\}$, of $B$ which covers $B$, as well as a collection of clopen collections $\mathcal{U}^\alpha$ whose union (over $\alpha$) is a cover of $f^{-1}(B)$.  By compactness of $f^{-1}(B)$, there is a finite collection $\alpha_1, \ldots, \alpha_n$ such that the union over the collections $\mathcal{U}^{\alpha_i}$ is a cover of $f^{-1}(B)$
%
%
%
%$\{G_{\alpha_j}\}_{j=1}^n$.  We want to make this cover disjoint in a particular way: first remove from each $G_{\alpha_j}$ $(j>1)$ any overlap with $G_{\alpha_1}$, then do the same for $G_{\alpha_2}$ and so on.  Then remove from $U_i^{\alpha_i}$ the (clopen) overlap with the pullback 
%
%need to know still a cover...




%Let $\{C_\alpha\}$ be the collection of $T$-components of $B$.  For each $\alpha$, there is a $(k, F, \mathcal{D}_f(F, S))$-cover $\mathcal{U}^\alpha = \{U_1^\alpha, \ldots, U^\alpha_k\}$ of $f^{-1}(C_\alpha)$ where each point is covered by at least $k-m$ sets.  Consider the cover $\mathcal{U} = \{\bigcup_\alpha U_1^\alpha, \ldots, \bigcup_\alpha U_k^\alpha\}$.  

\bibliography{mybibliography2.bib}
\bibliographystyle{plain}

\end{document}
% \bib, bibdiv, biblist are defined by the amsrefs package.
\begin{bibdiv}
\begin{biblist}

\bib{bartels}{incollection}{
      author={Bartels, Arthur},
       title={On proofs of the {F}arrell-{J}ones conjecture},
        date={2016},
   booktitle={Topology and geometric group theory},
      series={Springer Proc. Math. Stat.},
      volume={184},
   publisher={Springer, [Cham]},
       pages={1\ndash 31},
         url={https://doi.org/10.1007/978-3-319-43674-6_1},
      review={\MR{3598160}},
}

\bib{bartels_et_al}{article}{
      author={Bartels, Arthur},
      author={L\"{u}ck, Wolfgang},
      author={Reich, Holger},
       title={Equivariant covers for hyperbolic groups},
        date={200610},
     journal={Geom. Topol.},
      volume={12},
}

\bib{MR1808331}{article}{
      author={Bell, G.},
      author={Dranishnikov, A.},
       title={On asymptotic dimension of groups},
        date={2001},
     journal={Algebr. Geom. Topol.},
      volume={1},
       pages={57\ndash 71},
}

\bib{asymptotichurewicztheorem}{article}{
      author={Bell, Gregory},
      author={Dranishnikov, Alexander},
       title={A {H}urewicz-type theorem for asymptotic dimension and
  applications to geometric group theory},
        date={200408},
     journal={Transactions of the American Mathematical Society},
      volume={358},
}

\bib{asdimBD}{article}{
      author={Bell, Gregory},
      author={Dranishnikov, Alexander},
       title={Asymptotic dimension},
        date={200806},
     journal={Topology and its Applications},
      volume={155},
}

\bib{bonicke2021dynamic}{misc}{
      author={B{\"o}nicke, Christian},
      author={Dell'Aiera, Cl{\'e}ment},
      author={Gabe, James},
      author={Willett, Rufus},
       title={Dynamic asymptotic dimension and matui's {HK} conjecture},
        date={2021},
}

\bib{Brodskiy2006AHT}{article}{
      author={Brodskiy, Nikolay},
      author={Dydak, Jerzy},
      author={Levin, Michael},
      author={Mitra, Atish},
       title={A {H}urewicz theorem for the {A}ssouad--{N}agata dimension},
        date={2006},
     journal={J. London Math. Soc.},
      volume={77},
       pages={741\ndash 756},
}

\bib{https://doi.org/10.1112/blms/bdt045}{article}{
      author={Chen, Xiaoman},
      author={Wang, Qin},
      author={Wang, Xianjin},
       title={Characterization of the haagerup property by fibred coarse
  embedding into hilbert space},
        date={2013},
     journal={Bulletin of the London Mathematical Society},
      volume={45},
      number={5},
       pages={1091\ndash 1099},
  eprint={https://londmathsoc.onlinelibrary.wiley.com/doi/pdf/10.1112/blms/bdt045},
  url={https://londmathsoc.onlinelibrary.wiley.com/doi/abs/10.1112/blms/bdt045},
}

\bib{MR4091056}{article}{
      author={Delabie, Thiebout},
      author={Khukhro, Ana},
       title={Coarse fundamental groups and box spaces},
        date={2020},
     journal={Proc. Roy. Soc. Edinburgh Sect. A},
      volume={150},
      number={3},
       pages={1139\ndash 1154},
}

\bib{boxspacesDT}{article}{
      author={Delabie, Thiebout},
      author={Tointon, Matthew},
       title={The asymptotic dimension of box spaces of virtually nilpotent
  groups},
        date={2018},
     journal={Discrete Math.},
      volume={341},
       pages={1036\ndash 1040},
}

\bib{DELGADO2011398}{article}{
      author={Delgado, Alberto~L.},
      author={Robinson, Derek~J.S.},
      author={Timm, Mathew},
       title={Generalized baumslag--solitar groups and geometric
  homomorphisms},
        date={2011},
        ISSN={0022-4049},
     journal={Journal of Pure and Applied Algebra},
      volume={215},
      number={4},
       pages={398\ndash 410},
  url={https://www.sciencedirect.com/science/article/pii/S002240491000099X},
}

\bib{MR2213160}{article}{
      author={Dranishnikov, Alexander},
      author={Smith, J.},
       title={Asymptotic dimension of discrete groups},
        date={2006},
     journal={Fund. Math.},
      volume={189},
      number={1},
       pages={27\ndash 34},
}

\bib{engelking}{book}{
      author={Engelking, Ryszard},
       title={Dimension theory},
      series={North-Holland mathematical library},
   publisher={PWN},
        date={1978},
}

\bib{Exel2015PartialDS}{article}{
      author={Exel, Ruy},
       title={Partial dynamical systems, {F}ell bundles and applications},
        date={2015},
     journal={arXiv: Operator Algebras},
}

\bib{MR4066584}{article}{
      author={finiteness, Almost},
      author={the small~boundary property},
       title={Kerr, david and szab\'{o}, g\'{a}bor},
        date={2020},
     journal={Comm. Math. Phys.},
      volume={374},
      number={1},
       pages={1\ndash 31},
}

\bib{finnsell2015asymptotic}{misc}{
      author={Finn-Sell, Martin},
      author={Wu, Jianchao},
       title={The asymptotic dimension of box spaces for elementary amenable
  groups},
        date={2015},
}

\bib{asdimnote}{article}{
      author={Fujiwara, Koji},
      author={Whyte, Kevin},
       title={A note on spaces of asymptotic dimension one},
        date={2007},
     journal={Algebraic {\&} Geometric Topology},
      volume={7},
       pages={1063\ndash 1070},
}

\bib{https://doi.org/10.48550/arxiv.2201.03409}{misc}{
      author={Gardella, Eusebio},
      author={Geffen, Shirly},
      author={Kranz, Julian},
      author={Naryshkin, Petr},
       title={Classifiability of crossed products by nonamenable groups},
         how={arXiv},
        date={2022},
}

\bib{MR2431020}{article}{
      author={Gentimis, Thanos},
       title={Asymptotic dimension of finitely presented groups},
        date={2008},
        ISSN={0002-9939},
     journal={Proc. Amer. Math. Soc.},
      volume={136},
      number={12},
       pages={4103\ndash 4110},
         url={https://doi.org/10.1090/S0002-9939-08-08973-9},
      review={\MR{2431020}},
}

\bib{Gromov1991AsymptoticIO}{article}{
      author={Gromov, Mikhael},
       title={Asymptotic invariants of infinite groups in geometric group
  theory},
        date={1991},
     journal={Geometric group theory},
      volume={2},
}

\bib{MR87652}{article}{
      author={Gruenberg, Karl},
       title={Residual properties of infinite soluble groups},
        date={1957},
     journal={Proc. London Math. Soc. (3)},
      volume={7},
       pages={29\ndash 62},
}

\bib{Guentner2016DynamicAD}{article}{
      author={Guentner, E.},
      author={Willett, R.},
      author={Yu, Guoliang},
       title={Dynamic asymptotic dimension and controlled operator {K}-theory},
        date={2016},
     journal={arXiv: {K}-Theory and Homology},
}

\bib{dasdimGWY}{article}{
      author={Guentner, Erik},
      author={Willett, Rufus},
      author={Yu, Guoliang},
       title={Dynamic asymptotic dimension: relation to dynamics, topology,
  coarse geometry, and ${C^*}$-algebras},
        date={201510},
     journal={Math. Ann.},
      volume={367},
}

\bib{MR110750}{article}{
      author={Hall, P.},
       title={On the finiteness of certain soluble groups},
        date={1959},
     journal={Proc. London Math. Soc. (3)},
      volume={9},
       pages={595\ndash 622},
}

\bib{MR1143191}{article}{
      author={Hillman, Jonathan},
      author={Linnell, Peter},
       title={Elementary amenable groups of finite {H}irsch length are
  locally-finite by virtually-solvable},
        date={1992},
     journal={J. Austral. Math. Soc. Ser. A},
      volume={52},
      number={2},
       pages={237\ndash 241},
}

\bib{hillman_1991}{article}{
      author={Hillman, Jonathan~A.},
       title={Elementary amenable groups and 4-manifolds with {E}uler
  characteristic 0},
        date={1991},
     journal={J. Aust. Math. Soc.},
      volume={50},
      number={1},
       pages={160\ndash 170},
}

\bib{https://doi.org/10.48550/arxiv.2107.05273}{misc}{
      author={Kerr, David},
      author={Naryshkin, Petr},
       title={Elementary amenability and almost finiteness},
         how={arXiv},
        date={2021},
}

\bib{LiDynamics}{article}{
      author={Li, Xin},
       title={Dynamic characterizations of quasi-isometry, and applications to
  cohomology},
        date={2016},
     journal={Algebraic {\&} Geometric Topology},
      volume={18},
       pages={3477\ndash 3535},
}

\bib{MR285589}{article}{
      author={Meskin, Stephen},
       title={Nonresidually finite one-relator groups},
        date={1972},
     journal={Trans. Amer. Math. Soc.},
      volume={164},
       pages={105\ndash 114},
}

\bib{munkres}{book}{
      author={Munkres, James},
       title={Topology},
   publisher={Prentice Hall},
        date={2000},
}

\bib{Roe2003LecturesOC}{book}{
      author={Roe, John},
       title={Lectures on coarse geometry},
      series={University Lecture Series},
   publisher={American Mathematical Society, Providence, RI},
        date={2003},
      volume={31},
        ISBN={0-8218-3332-4},
         url={https://doi.org/10.1090/ulect/031},
      review={\MR{2007488}},
}

\bib{warpedcones}{article}{
      author={Sawicki, Damian},
      author={Kielak, Dawid},
       title={Warped cones, (non-)rigidity, and piecewise properties},
        date={201809},
     journal={Proc. London Math. Soc.},
      volume={118},
}

\bib{MR223439}{article}{
      author={Stallings, John},
       title={Groups of dimension 1 are locally free},
        date={1968},
        ISSN={0002-9904},
     journal={Bull. Amer. Math. Soc.},
      volume={74},
       pages={361\ndash 364},
         url={https://doi.org/10.1090/S0002-9904-1968-11955-X},
      review={\MR{223439}},
}

\bib{szabo2017rokhlin}{article}{
      author={Szab\'o, G\'abor},
      author={Wu, Jianchao},
      author={Zacharias, Joachim},
       title={Rokhlin dimension for actions of residually finite groups},
        date={2019},
     journal={Ergodic Theory and Dynamical Systems},
      volume={39},
      number={8},
       pages={2248\ndash 2304},
}

\bib{exponentialgrowthBSgroups}{article}{
      author={ucher Karlsson, Michelle},
      author={Talambutsa, Alexey},
       title={Minimal exponential growth rates of metabelian
  {B}aumslag-{S}olitar groups and lamplighter groups},
        date={201506},
     journal={Groups, Geometry, and Dynamics},
      volume={11},
}

\bib{notes_on_prop_A}{article}{
      author={Willett, Rufus},
       title={Some notes on property {A}},
        date={200701},
     journal={Limits of Graphs in Group Theory and Computer Science},
}

\end{biblist}
\end{bibdiv}
